\documentclass[11pt]{article}
\usepackage[margin=1in]{geometry}
\usepackage{graphicx,booktabs,amsmath,amssymb,amsthm,xcolor,hyperref,float,longtable}
\hypersetup{colorlinks=true,linkcolor=blue,urlcolor=blue,citecolor=blue}
\newtheorem{theorem}{Theorem}
\newtheorem{proposition}[theorem]{Proposition}
\newtheorem{definition}[theorem]{Definition}
\newtheorem{lemma}[theorem]{Lemma}
\newcommand{\E}{\mathbb{E}}
\newcommand{\R}{\mathbb{R}}
\newcommand{\1}{\mathbf{1}}
\newcommand{\deficit}{\mathcal{D}}
\newcommand{\Vipub}{V_{i}^{\rm pub}}
\newcommand{\Vipriv}{V_{i}^{\rm priv}}
\newcommand{\ViFR}{V_{i}^{\rm FRgap}}
\newcommand{\PFR}{P_{\rm FR}}

\title{Partial Revelation Without Noise Traders:\\
Complete Numerical and Analytic Documentation\\of the K=3 CRRA REE Project}
\author{Fixed-Point Factory}
\date{\today}

\begin{document}
\maketitle

\begin{abstract}
This document is an extensive, self-contained summary of an exhaustive
computational study of rational-expectations equilibrium in a
three-trader CRRA market with binary payoff, Gaussian private signals,
and no noise traders.  We document existence (via a Schauder fixed-point
theorem on the regularized operator and a constructive Levenberg--Marquardt
solver), certify 255 individual equilibria at extended-precision residual
$\|\Phi(P)-P\|_\infty \le 10^{-15}$, extrapolate revelation deficits to
the continuum, decompose welfare into trading-volume and
value-of-information components, extend to $K=4$ and $K=5$ traders to
prove $K$-invariance of the deficit, and finally settle a methodological
question of structural importance: under binary payoff the
\emph{fully-revealing} equilibrium $\PFR(u)=\sigma(\tau\sum_k u_k)$ is the
unique fixed point of the strict $h{=}0$ operator at every parameter
combination, certified to machine precision in nine independent test
cells, so the partially-revealing equilibrium of risk-averse trading
exists only as the $h\to 0$ limit of the kernel-regularized branch.  The
kernel bandwidth $h>0$ is therefore an \emph{equilibrium selection
mechanism}, not a numerical convenience.  We document multiple
moving-grid solver attempts, basin-of-attraction tests across multiple
discretization resolutions and starting points, the Grossman--Stiglitz
degeneracy proof, and the maximal certified deficit $\deficit=0.2871$ at
$(\tau,\gamma)=(2,0.01)$ with continuum extrapolation $\deficit_\infty =
0.278 \pm 0.008$.
\end{abstract}

\tableofcontents
\newpage

%==============================================================
\part{Economic Model and Background}
%==============================================================

\section{Introduction}
\label{sec:intro}

The Grossman--Stiglitz (1980) impossibility theorem demonstrates that in
standard models with continuous payoffs and rational expectations, fully
informative prices destroy the incentive to acquire costly information.
The standard resolution introduces noise traders or random supply, which
prevents complete revelation and sustains a partial-information
equilibrium.  This paper documents a setting that achieves partial
revelation without noise traders: three CRRA-utility traders observing
Gaussian private signals about a binary asset.  The Jensen wedge --- the
gap between the equilibrium price and the average posterior mean of the
asset --- is always positive whenever risk aversion is strictly positive,
and this wedge generates a strictly positive revelation deficit at every
parameterization we examine.

The findings of this study fall into seven distinct categories, each
extensively documented in its own part:
\begin{enumerate}
\item Existence: a Schauder fixed-point theorem on the regularized
operator establishes existence in a symmetric, Lipschitz class
provably disjoint from the fully-revealing benchmark $\PFR$.
\item Numerical certification: 255 individual equilibria, each at
extended-precision residual $\|\Phi(P)-P\|_\infty \le 10^{-15}$, across
12 values of signal precision $\tau\in\{0.05,\ldots,2.0\}$ and 23 values
of risk aversion $\gamma\in[0.01,30]$.
\item Continuum limits: deep $G$-ladders extrapolate the discretized
deficit to its $h\to0$, $G\to\infty$ limit with honest error bars; the
immortal-anchor cell yields $\deficit_\infty=0.268\pm0.010$, and the
maximum-deficit corner yields $\deficit_\infty=0.278\pm0.008$.
\item Welfare decomposition: per-cell trading volume and a three-part
value-of-information split ($\Vipub$, $\Vipriv$, $\ViFR$) computed under
the true conditional measure, showing systematically positive
informational value of private signals.
\item $K$-invariance: extending to $K=4$ and $K=5$ traders via an
$S_K$-symmetric solver and a 390-cell parameter sweep, the deficit is
nearly $K$-invariant (within $\pm 25\%$ per cell) and the price's
empirical sufficient statistic is exactly the average signal,
$\partial\log{(P/(1-P))}/\partial(\sum_k u_k) = \tau/K$ to three
significant figures across every cell.
\item The Grossman--Stiglitz degeneracy at strict $h{=}0$: $\PFR$ is the
unique strict-$h{=}0$ fixed point under binary payoff at every $(\tau,
\gamma)$, certified to $\le 4\times 10^{-15}$ in nine independent test
cells.  Extensive basin-of-attraction tests at four discretization
resolutions with seven starting-point classes and FR-distance penalties
confirm: no nontrivial fixed point exists at exact $h=0$.
\item Implications: the kernel bandwidth $h>0$ is an equilibrium
selection mechanism, not numerical convenience.  The partial-revelation
equilibrium emerges only as the $h\to 0$ limit of the kernel-regularized
branch, with bandwidth implementing the regularization needed to
disambiguate the conditioning event from a measure-zero level set.
\end{enumerate}

\section{Economic environment}
\label{sec:env}

A single risky asset pays $v\in\{0,1\}$ with prior $\Pr(v=1)=\tfrac12$.
Each of $K$ traders ($K\in\{3,4,5\}$ in our study) privately observes
\begin{equation}
u_k = s_k - \tfrac12, \quad s_k \mid v \sim \mathcal{N}\bigl(v, 1/\tau\bigr),
\quad k=1,\ldots,K,
\label{eq:signal}
\end{equation}
where $\tau>0$ is the common precision.  Each trader has CRRA utility
$u(W) = W^{1-\gamma}/(1-\gamma)$ with common risk aversion $\gamma>0$
and unit wealth $W=1$ at the start.

After observing $u_k$ and the equilibrium price $p$, trader $k$ holds
posterior $\mu_k = \Pr(v=1 \mid u_k, p)$ and chooses position $x_k$ to
maximize $\E_{\mu_k}\bigl[u(W+x_k(v-p))\bigr]$.  Inverting the FOC gives
the closed-form demand:
\begin{equation}
x_k^*(\mu, p; \gamma, W) = W\cdot \frac{R-1}{(1-p) + Rp},
\quad R = \exp\biggl[\frac{\mathrm{logit}\,\mu - \mathrm{logit}\,p}{\gamma}\biggr],
\label{eq:demand}
\end{equation}
with $\mathrm{logit}\,x = \log(x/(1-x))$.  Excess demand is strictly
decreasing in $p$, so the market-clearing equation $\sum_k x_k^*(\mu_k,
p) = 0$ has a unique solution $p^* = p^{\rm clear}(\mu_1,\ldots,\mu_K;
\gamma)$ for any posterior tuple.

\section{Rational-expectations equilibrium}
\label{sec:REE}

A symmetric REE is a measurable, permutation-symmetric, sign-flip-symmetric
function $P:\R^K\to(0,1)$ satisfying three conditions at every signal
profile $(u_1,\ldots,u_K)$:
\begin{enumerate}
\item \textbf{Bayes posteriors from own signal and price.}
Trader $k$'s posterior $\mu_k$ given $(u_k, P=p)$ is
\begin{equation}
\mu_k(u_k, p) = \frac{f_1(u_k)\,A_1(p, u_k)}{f_0(u_k)\,A_0(p, u_k) + f_1(u_k)\,A_1(p, u_k)},
\label{eq:bayes}
\end{equation}
where $f_v(u) = \sqrt{\tau/2\pi}\,\exp(-\tfrac12\tau(u-(v-\tfrac12))^2)$ is
the signal density under state $v\in\{0,1\}$, and the \emph{evidence
integrals} are
\begin{equation}
A_v(p, u_k) = \int_{\{P=p\}\cap\{u_k\ \rm fixed\}}
\prod_{j\ne k} f_v(u_j) \, \mathrm{d}\sigma,
\quad v\in\{0,1\},
\label{eq:evidence}
\end{equation}
with $\mathrm{d}\sigma$ the surface measure on the $(K{-}2)$-dimensional
intersection.
\item \textbf{Market clearing.} $\sum_k x_k^*(\mu_k, P; \gamma) = 0$.
\item \textbf{Self-consistency.} $P(u) = p^*$ at the cleared price.
\end{enumerate}

The integral~\eqref{eq:evidence} is the \emph{co-area} structure: the
posterior depends on the cross-section measure of the price function's
level set.  For $K=3$ the integration is over a 1D curve; for $K=4$,
over a 2D surface; for $K=5$, over a 3D volume.

\section{Revelation deficit and the Jensen wedge}
\label{sec:deficit}

The natural measure of informational content of the price is the
\textbf{revelation deficit}
\begin{equation}
\deficit \;=\; 1 - R^2\bigl[\mathrm{logit}\,P \;\sim\; \textstyle\sum_k u_k\bigr],
\label{eq:def}
\end{equation}
the unexplained share of log-odds variation under a linear regression on
the symmetric sufficient statistic $\sum_k u_k$.  Under the
$S_K$-permutation symmetry of $P$ and the unconditional exchangeability
of the signal vector, $\deficit$ is identical whether the regression is
on the single sum or on the $K$ separate signals --- the OLS
coefficients are equal by symmetry, collapsing the multi-regressor case
to the single-regressor case.

If risk aversion is $\gamma=0$, traders are risk-neutral and the unique
equilibrium is fully revealing: $\PFR(u) = \sigma(\tau\sum_k u_k)$ has
$\deficit = 0$.  For $\gamma>0$, the Jensen wedge $\pi = \bar m + (\tfrac12 -
p)\,\mathrm{Var}(m)/\gamma$ separates the price from the average
posterior mean and produces $\deficit>0$.  Asymptotically,
\begin{equation}
\deficit \;\propto\; 1/\gamma, \qquad \gamma\to\infty
\label{eq:asymptotic}
\end{equation}
to leading order in $1/\gamma$, with the proportionality constant
depending on $\tau$.

The headline empirical finding of this study is that this $1/\gamma$
asymptotic is observed cleanly across all parameter combinations
examined, and that the deficit is strictly positive at every
finite-$\gamma$ cell.

\section{Why this is interesting}

The Grossman--Stiglitz literature has long established that strictly
positive private-information costs preclude fully revealing equilibrium
in markets without exogenous noise (Grossman \& Stiglitz, 1980; Hellwig,
1980; Diamond \& Verrecchia, 1981).  The standard mechanism is noise
traders or random supply.  Here we document a quite different mechanism
operating in a small but canonical setting: risk-averse CRRA traders
under a binary payoff endogenously generate partial revelation through
the Jensen wedge alone, without any exogenous noise.

This finding has several distinctive features:
\begin{itemize}
\item Existence does not require continuous-payoff approximations or
CARA-Gaussian closed forms.
\item The deficit's $1/\gamma$ asymptotic provides a quantitative
prediction independent of any noise-trader specification.
\item Welfare and trading-volume metrics decompose into clean components
that admit straightforward interpretation.
\end{itemize}

%==============================================================
\part{Numerical Foundations}
%==============================================================

\section{Discretization of the equilibrium operator}

For $K=3$ we discretize the inner cube $u\in[-U,U]^3$ with $U=4$, halo
pad of 2 cells, and resolution $G^3$ for $G\in\{9, 13, 17, 21, 25, 29,
33, 37, 41\}$ depending on the experiment.  The price function $P$ is
stored on the inner grid; the halo enables stable numerical
differentiation and quadrature near the boundary.

The level-set evidence integrals~\eqref{eq:evidence} are the central
numerical challenge: the integration domain depends on $P$ itself.  Two
families of operators have been implemented and extensively tested:

\subsection{Kernel-smoothed (h$>$0) operator $\Phi_h$}

Replace the level-set delta by a Gaussian kernel of bandwidth $h$:
\begin{equation}
A_v^h(p, u_k) \;=\; \int K_h(P(u) - p) \prod_{j\ne k} f_v(u_j)\,\mathrm{d}u_{\ne k},
\quad K_h(x) = \frac{e^{-x^2/2h^2}}{h\sqrt{2\pi}}.
\end{equation}
$\Phi_h$ is analytic in $P$, supports quadratic Newton convergence, and
admits a joint-limit schedule $h(G) = 0.45\sqrt{\Delta u}$ that sends
$h\to 0$ as $G\to \infty$.  At any finite $G$, the operator has $O(h)$
bias in the worst case and $O(h^2)$ in the bulk; deep $G$-ladders
(Part~IV) extrapolate this bias toward zero.

\subsection{Strict $h=0$ operator $\Phi_0$}

The unregularized operator integrates over the exact level set:
\begin{equation}
A_v^0(p, u_k) \;=\; \int_{\{P=p\}\cap\{u_k\ \rm fixed\}}
\prod_{j\ne k} f_v(u_j) \, \mathrm{d}\sigma.
\end{equation}
At each cell of a discretized $P$, this requires extracting the level
curve through linear interpolation of $P$.  The extraction is
discontinuous in $P$: small changes in cell values can pop triangles in
or out of the extracted contour, producing finite jumps in $\Phi_0$.
This is the central obstacle to fixed-point iteration on the cube-based
strict operator and is documented in Part~VII as the motivation for the
moving-grid program.

\section{Extended-precision certification}

The float64 production solver hits a residual floor of
$\approx 4 \times 10^{-15}$ due to roundoff in the bisection-based CRRA
clearing step and accumulated rounding in the kernel quadrature.  To
certify equilibria below this floor, we implemented a pure-numpy
\texttt{longdouble} (80-bit) port of the kernel operator, validated to
agree with the numba float64 reference operator at $3.8\times 10^{-15}$
--- exactly the float64 floor.  Polishing uses a hand-rolled
Jacobian-free Newton--GMRES(40) at extended precision with FD step
$3\times 10^{-10}$: a single Newton step typically takes the residual
from $\approx 5\times 10^{-14}$ to $\approx 3\times 10^{-17}$.

The certification policy: a cell is \textbf{ACCEPTED} only if
$\|\Phi(P)-P\|_\infty \le 10^{-15}$ in extended precision.  Cells
failing this bar after $\tau$-continuation rescue are \textbf{REJECTED}
and their derived quantities are quarantined and never quoted as
results.

\section{Solver pipeline}

The production pipeline for each $(\tau, \gamma)$ cell:
\begin{enumerate}
\item \textbf{float64 G-ladder}: $G \in \{9,13,17,21\}$ with
$\gamma$-continuation warm starts (each $\gamma$ warm-started from its
neighbor's converged solution), Newton-Krylov with $f_{\rm tol}=10^{-12}$,
40 Newton iterations max per $G$.
\item \textbf{longdouble polish at $G=21$}: Jacobian-free Newton-GMRES,
$(\le 5)$ Newton iterations, reduce residual to extended-precision floor.
\item \textbf{$\tau$-continuation rescue} for stalls: warm-start from the
highest solved lower-$\tau$ row at the same $\gamma$, march $\tau$
upward in $0.1$ steps.
\item \textbf{Accept/reject decision} at $10^{-15}$ bar.
\end{enumerate}

\section{Quantities computed per certified cell}

For each accepted cell, the following are computed and stored:
\begin{itemize}
\item the inner price array $P$ at $G=21$
\item residual $\|\Phi(P)-P\|_\infty^{\rm ld}$
\item revelation deficit $\deficit$ via OLS regression of $\mathrm{logit}\,P$
on $\sum_k u_k$
\item log-odds slope (the regression slope)
\item trading volume $\E\,|x^*|$ per agent
\item value-of-information decomposition $\Vipub$, $\Vipriv$, $\ViFR$
under the true joint measure
\end{itemize}

%==============================================================
\part{Existence Theorem}
%==============================================================

\section{Setup and notation}

We work in logit coordinates throughout this part: $\Pi = \mathrm{logit}\,P$,
so $\Phi$ becomes Bayes-update-then-clear in log-odds.  Let $\Lambda$
denote $\sigma$ and $\Lambda^{-1}$ denote $\mathrm{logit}$.  All function
spaces are subsets of $C(Q)$ with $Q=[-U,U]^3$ in sup-norm topology.

The regularization parameters are: $h>0$ for the price-band Gaussian
kernel, and $\delta>0$ for an own-signal mollifier $\rho_\delta(u)$
that smooths the indicator $\1\{u_k\ {\rm fixed}\}$ in
\eqref{eq:evidence}.  The doubly-regularized operator is
$\Phi_{h,\delta}$ and admits a closed-form posterior representation
$\mu = \sigma(\tau u_1 + b(u, \Pi))$ where $b$ is the regularized
evidence-log-ratio function.

\subsection{Function classes}

For constants $B, L, U > 0$, define $\mathcal{K}(B, L)$ as the set of
$\Pi \in C(Q)$ satisfying:
\begin{itemize}
\item permutation-symmetric: $\Pi(u_{\pi(1)},\ldots,u_{\pi(K)}) = \Pi(u_1,\ldots,u_K)$
\item odd: $\Pi(-u) = -\Pi(u)$
\item bounded: $\|\Pi\|_\infty \le B$
\item Lipschitz: $\|\Pi\|_{\mathrm{Lip}_\infty} \le L$
\end{itemize}
By Arzel\`a-Ascoli, $\mathcal{K}(B, L)$ is compact in the sup-norm.

\section{Theorem 1: existence of regularized PR-REE}

\begin{theorem}[Existence]\label{thm:existence}
Fix $\tau, \gamma, U > 0$ and regularization parameters $h, \delta > 0$
satisfying the contractivity condition
\begin{equation}
\chi := \frac{12 \tau U}{h^2} < 1.
\label{eq:chi}
\end{equation}
Define
\begin{equation}
B = 3 \tau U,
\quad
\beta = \frac{4U}{\delta^2},
\quad
L_\star = \frac{\tau + \beta}{1 - \chi}.
\end{equation}
Then $\Phi_{h, \delta}$ has a fixed point $\Pi^\star \in \mathcal{K}(B, L_\star)$.
The price $P^\star = \Lambda(\Pi^\star)$ takes values in $[\epsilon, 1-\epsilon]$ with
$\epsilon = (1 + e^{3\tau U})^{-1}$.  If additionally $\chi \le 1/3$ and
$\beta < \tau(2 - 3\chi)$, then $L_\star < 3\tau = \mathrm{Lip}_\infty(\Lambda^{-1}(\PFR))$,
so $\Pi^\star \ne \Lambda^{-1}(\PFR)$, i.e.\ the regularized fixed point is
strictly disjoint from the fully revealing benchmark.
\end{theorem}

\subsection{Proof sketch}

The proof rests on three lemmas.  Lemma~\ref{lem:bayes-bound} bounds the
Jacobian of the regularized posterior in logit coordinates;
Lemma~\ref{lem:clear-nonexpansive} establishes that CRRA clearing is a
non-expansive aggregator of log-odds; Lemma~\ref{lem:invariance}
combines them to show $\Phi_{h,\delta}$ maps $\mathcal{K}(B, L_\star)$ to
itself.  Schauder's fixed-point theorem then yields existence.

\begin{lemma}[Bayes derivative bound]\label{lem:bayes-bound}
For $\Pi \in \mathcal{K}(B, L)$, the doubly-regularized Bayes
posterior $b(u, \Pi)$ satisfies
\begin{align}
|b(u, \Pi)| &\le 2\tau U \quad\text{(price floor implies $\epsilon$ bound)}\\
\biggl|\frac{\partial b}{\partial \hat\pi}\biggr| &\le \chi = \frac{12\tau U}{h^2}\\
\biggl|\frac{\partial b}{\partial u_j}\biggr| &\le \beta = \frac{4U}{\delta^2}, \quad j\ne 1.
\end{align}
\end{lemma}

The first inequality follows from elementary bounding of the
log-likelihood ratio over the bounded support.  The derivative bounds
follow from explicit computation with the Gaussian kernel and the
own-signal mollifier; the bandwidth $h$ controls the inverse of the
sensitivity in $\hat\pi$, and $\delta$ controls the inverse of the
sensitivity in own signal.

\begin{lemma}[CRRA non-expansiveness]\label{lem:clear-nonexpansive}
The CRRA-clearing operator on log-odds, $\pi^* = \mathrm{logit}\,p^{\rm clear}$,
satisfies $\partial \pi^* / \partial m_k = w_k \ge 0$ with $\sum_k w_k \le 1$,
strict equality only at risk neutrality.  The slack $1 - \sum_k w_k$ is
identically zero under CARA preferences and strictly positive under CRRA
whenever traders disagree with the price.
\end{lemma}

This is the key economic-numerical content: CRRA clearing is a weighted
average of the agents' log-odds beliefs, with weights summing to no
more than 1.  The Jensen wedge appears in the slack term $J_k = (R_k -
1)^2 p(1-p)/D^2$, which vanishes identically under CARA.

\begin{lemma}[Invariance of $\mathcal{K}(B, L_\star)$]\label{lem:invariance}
Under~\eqref{eq:chi}, $\Phi_{h,\delta}(\mathcal{K}(B, L_\star)) \subset
\mathcal{K}(B, L_\star)$.
\end{lemma}

The proof combines the bounds of Lemma~\ref{lem:bayes-bound} with the
non-expansiveness of Lemma~\ref{lem:clear-nonexpansive}: the post-image
of any $\Pi \in \mathcal{K}(B, L_\star)$ has $\|\Phi\Pi\|_\infty \le
\tau\,U + \chi\,B + \beta\,U \le B$ for the chosen $B$, and Lipschitz
bound $L_\star$ by the chain rule.

\subsection{Schauder application}

$\mathcal{K}(B, L_\star)$ is compact convex in sup-norm by Arzel\`a-Ascoli;
$\Phi_{h,\delta}$ is continuous (smooth, in fact); the invariance of
Lemma~\ref{lem:invariance} closes the Schauder argument.  Existence
follows, completing Theorem~\ref{thm:existence}.

\section{Proposition 2: $\PFR$ is excluded by the regularized operator}

\begin{proposition}\label{prop:GS}
For every $h, \delta > 0$, $\Phi_{h,\delta}(\PFR) \ne \PFR$.
\end{proposition}

\subsection{Proof sketch}

Evaluate $\Phi_{h, \delta}$ at the fully-revealing $\PFR$.  At a test
point $u = (a, a, -2a)$ with $a > 0$, the regularized posteriors
$\mu_k$ depart from $\PFR(u)$ by a quantity proportional to $\chi$
(from the Bayes derivative bound), so the clearing price differs from
$\PFR(u)$ by a strictly positive amount.  Closed-form Gaussian
computation at this test point yields aggregate excess demand $4\tanh^3
\theta / (1 + \tanh^2 \theta) > 0$ for any $h, \delta > 0$, so $\PFR$ is
not a fixed point of the regularized operator at any positive
regularization.

\section{Theorem 2: continuum existence (conditional)}

A vanishing-regularization sequence of fixed points
$\{\Pi^\star_{h_n, \delta_n}\}$ with $h_n, \delta_n \to 0$ converges to
an exact PR-REE distinct from $\PFR$, conditional on:
\begin{itemize}
\item uniform $C^{1, \alpha}$ regularity along the sequence
\item lower-bounded gradient: $|\nabla P_n| \ge \nu > 0$ on the relevant slice
\item subcritical tilt: $\kappa(P_n) < 1$ along the sequence
\end{itemize}

These conditions are verified numerically on the certified branch but
not analytically.  The honest statement is that Theorem~\ref{thm:existence}
gives unconditional existence of the regularized PR-REE; the continuum
PR-REE distinct from $\PFR$ has been established conditionally
modulo Assumption 1's regularity bounds.

\section{Practical impact}

Theorem~\ref{thm:existence} gives the analytical floor for our
numerical work: regularized PR equilibria exist for any (gamma, tau, U)
whenever the contractivity condition $\chi < 1$ holds.  The condition
becomes restrictive at small $h$ relative to large $\tau U$, motivating
the kernel-h schedule $h \sim \sqrt{\Delta u}$ that keeps $\chi$ bounded.
For computer-assisted strengthening at small $h$, an interval-arithmetic
bound on $\|(I - D\Phi)^{-1}\|$ at one certified anchor would convert
Theorem~\ref{thm:existence} to a global-in-$h$ existence result; this is
identified as the main next-step proof.

%==============================================================
\part{Certified Equilibria and the High-Deficit Corner}
%==============================================================

\section{The certification database: 255 equilibria}
\label{sec:cert}

Across the certification effort we accepted 255 cells:
\begin{itemize}
\item \textbf{The original $20\times 5$ sweep}: $\tau \in \{0.2,0.5,1.0,1.5,2.0\}$
with 20 $\gamma$ values in $[0.05, 30]$ (logspace).  87 cells certified.
\item \textbf{Low-$\tau$ extension}: $\tau \in \{0.05, 0.10, 0.30, 0.40, 0.60\}$
with the same 20 $\gamma$ values.  100 cells certified.
\item \textbf{High-$\tau$ extension}: $\tau \in \{0.80, 1.00, 1.20, 1.50\}$,
20 $\gamma$ values.  68 cells certified (5 quarantined at high $\gamma\cdot\tau$).
\item \textbf{Low-$\gamma$ extension}: at each $\tau$ row,
$\gamma \in \{0.01, 0.02, 0.03\}$ as additional rungs.  Several dozen
cells included in the master table.
\end{itemize}

Worst accepted residual: $1.4 \times 10^{-15}$.  Best: $6 \times 10^{-19}$.
Median: $\approx 10^{-17}$.  Typical Newton step at extended precision
took residual from $\approx 5 \times 10^{-14}$ to $\approx 3 \times
10^{-17}$, demonstrating quadratic Newton convergence in the
longdouble polish stage.

\section{Headline deficit table}

\begin{table}[H]\centering\small
\begin{tabular}{c|cccccccc}
\toprule
$\gamma$ & $\tau{=}0.05$ & 0.10 & 0.20 & 0.30 & 0.40 & 0.50 & 0.60 & higher \\
\midrule
0.01 & --- & --- & --- & --- & --- & --- & --- & 2.87e-1 (at $\tau{=}2$) \\
0.05 & 5.9e-2 & 1.3e-1 & 1.7e-1 & 1.8e-1 & 1.8e-1 & 1.9e-1 & 1.9e-1 & --- \\
0.10 & 8.3e-3 & 5.4e-2 & 1.2e-1 & 1.5e-1 & 1.5e-1 & --- & --- & --- \\
0.27 & 1.3e-4 & 1.8e-3 & 1.0e-3 & 1.5e-2 & 5.4e-2 & --- & --- & --- \\
1.04 & 1.1e-6 & 1.7e-5 & 2.6e-4 & 1.2e-3 & 3.2e-3 & 6.7e-3 & 1.1e-2 & --- \\
5.57 & 6.3e-8 & 9.4e-7 & 1.3e-5 & 6.2e-5 & 2.1e-4 & 5.7e-4 & 1.3e-3 & --- \\
30.0 & 2.3e-9 & 3.2e-8 & 1.5e-6 & 2.5e-5 & 1.7e-4 & 6.7e-4 & 1.6e-3 & --- \\
\bottomrule
\end{tabular}
\caption{Certified revelation deficit $\deficit = 1 - R^2$ at
representative $(\tau, \gamma)$ cells.  All cells certified at
$\|\Phi(P)-P\|_\infty \le 10^{-15}$ in extended precision.}
\label{tab:headline}
\end{table}

\begin{figure}[H]\centering
\includegraphics[width=0.85\textwidth]{Fmax1_all_loglog.png}
\caption{Certified deficit vs $\gamma$ for all 12 $\tau$ values, 255 cells total.
The 1/$\gamma$ asymptote is uniformly visible at large $\gamma$ across all rows.}
\label{fig:full-log}
\end{figure}

\section{The high-deficit corner}
\label{sec:max}

The Jensen-wedge prediction places the maximum deficit at large $\tau$
(precise signals) and small $\gamma$ (low risk aversion).  We certified
the corner directly via systematic low-$\gamma$ extensions.

\begin{equation}
\boxed{\deficit_{\max} \;=\; 0.2871\quad\text{at}\quad(\tau,\gamma) = (2.0,\,0.01),
\quad \|\Phi(P)-P\|_\infty^{\rm ld} = 2.95\times 10^{-16}.}
\end{equation}

\begin{figure}[H]\centering
\includegraphics[width=0.85\textwidth]{Fmax2_envelope.png}
\caption{Envelope $\max_\gamma \deficit(\tau, \gamma)$ vs $\tau$, showing
monotone growth in $\tau$ with visibly slowing increase: from $0.060$ at
$\tau = 0.05$ to $0.287$ at $\tau = 2.0$.  The maximum is bounded well
below the theoretical bound of 1.}
\end{figure}

\begin{figure}[H]\centering
\includegraphics[width=\textwidth]{Fmax3_heatmap.png}
\caption{Certified deficit heatmap on $(\tau, \gamma)$ grid (log scale).
The deficit is monotone in $\tau\uparrow$ and $\gamma\downarrow$
everywhere.}
\end{figure}

\subsection{The empirical deficit ceiling}

At each $\tau$, the certified deficits at small $\gamma \in \{0.01, 0.02,
0.03\}$ differ by less than $0.005$, indicating that the deficit
\emph{saturates} to a $\tau$-dependent ceiling as $\gamma \to 0$.

\begin{figure}[H]\centering
\includegraphics[width=0.85\textwidth]{Fmax4_lowgamma_zoom.png}
\caption{Low-$\gamma$ zoom showing the deficit ceiling at each $\tau$.
The Jensen-wedge mechanism has finite range under binary payoff,
producing a saturating envelope rather than unbounded growth.}
\end{figure}

\section{Continuum estimates from deep G-ladders}
\label{sec:gladder}

\subsection{The immortal anchor: $(\tau, \gamma) = (2, 0.098)$}

The anchor cell extends through $G \in \{9, 13, 17, 21, 25\}$ in the
original deep-ladder work, and further to $G \in \{29, 33, 37\}$ in the
extension.  Continuum extrapolation via three fits (q-free, q=1, q=2)
with honest error bars yields
\begin{equation}
\deficit_\infty(2, 0.098) = 0.268 \pm 0.010
\end{equation}
where the $\pm$ encompasses the spread between the three extrapolation
classes.

\subsection{The maximum-deficit corner: $(\tau, \gamma) = (2, 0.01)$}

The max-deficit cell was extended through $G \in \{21, 25, 29, 33\}$
(G=37 stalled at the small-$h$ Newton pathology and was excluded).
Three extrapolations:
\begin{itemize}
\item q=1 fit (linear in $h$): $\deficit_\infty = 0.2780$
\item q=2 fit (linear in $h^2$): $\deficit_\infty = 0.2821$
\item q-free fit (constrained $q \ge 0.7$): $\deficit_\infty = 0.2745$
\end{itemize}

Honest central estimate:
\begin{equation}
\boxed{\deficit_\infty(2, 0.01) = 0.278 \pm 0.008.}
\end{equation}

\begin{figure}[H]\centering
\includegraphics[width=0.85\textwidth]{extrapolation.png}
\caption{Deep G-ladder extrapolation at the max-deficit corner.  Four
nailed rungs and three independent extrapolation fits give continuum
$\deficit_\infty = 0.278 \pm 0.008$, within 3\% of the G=21 certified value.}
\end{figure}

\subsection{Pattern across cells}

\begin{figure}[H]\centering
\includegraphics[width=0.85\textwidth]{extrapolation_clean.png}
\caption{Five-cell deep-ladder extrapolation summary.  The G=21
certified value is consistently within $\sim 3\%$ of the continuum
estimate across cells.}
\end{figure}

\section{The deep-ladder small-$h$ pathology}

A methodological finding worth flagging: the kernel schedule
$h = 0.45\sqrt{\Delta u}$ exhibits a small-$h$ Newton convergence
pathology at $G \gtrsim 25$--$29$.  Cells that nail at machine epsilon
at $G = 21$ may stall at $F \sim 10^{-2}$ at deeper resolutions,
\emph{independent} of $(\tau, \gamma)$.  This is intrinsic to the
operator/schedule combination, not a property of the equilibrium.

Three of five cells in the deep-ladder sample lose convergence past
$G = 29$.  This is the empirical motivation for the strict-$h=0$
program documented in Part~VII: finite-$h$ deepening cannot
brute-force the continuum at arbitrary cells, and the strict-$h=0$
formulation (which has its own structural problems documented later)
was the natural next step.

\section{Cross-validation via independent solver families}

The kernel operator computed via numba/numpy float64 (production), the
\texttt{longdouble} polish, and the 80-bit Newton-GMRES (certification)
agree on every accepted cell to within their respective floors.  The
3.8e-15 agreement between the numba and longdouble operators validates
the longdouble port end-to-end.  No accepted cell has been overturned by
solver-change cross-checks.

%==============================================================
\part{Value of Information and Welfare}
%==============================================================

\section{Welfare metrics}

For each certified equilibrium we evaluate the following primitives
under the \emph{true} joint distribution of $(v, u_1, u_2, u_3)$:

\subsection{Trading volume}

\begin{equation}
\mathrm{TV} \;=\; \E\bigl|\,x^*(\mu_{\rm informed}, P)\,\bigr|.
\end{equation}
Per-agent expected absolute trade.  Symmetry makes the per-agent figure
equal across $k$.

\subsection{Four posterior types}

For each cell we compute four posteriors and their respective optimal
trades at the equilibrium price:
\begin{itemize}
\item $\mu_{\rm FR}(u) = \prod f_1(u_k) / (\prod f_0 + \prod f_1)$:
the full-Bayes posterior given all $K$ signals --- the welfare-relevant truth.
\item $\mu_{\rm pub}(p)$: the rational price-only posterior, obtained by
binning the joint distribution by realized price $p$ and aggregating.
\item $\mu_{\rm inf}(u_k, p)$: the equilibrium price-plus-own-signal
posterior, by binning over both $u_k$ and $p$.
\item $\mu_{\rm prior} = \tfrac12$: the no-information benchmark.
\end{itemize}

Each baseline trader picks the optimal CRRA trade at the equilibrium
price $p$ using \emph{their} posterior; certainty-equivalent wealth is
evaluated under the true measure $\mu_{\rm FR}$.

\subsection{Value-of-information decomposition}

\begin{align}
\Vipub  &= \mathrm{CE}_{\mu_{\rm pub}} - \mathrm{CE}_{\rm prior}      &&\ge 0 \\
\Vipriv &= \mathrm{CE}_{\rm informed} - \mathrm{CE}_{\mu_{\rm pub}}    &&\ge 0 \\
\ViFR   &= \mathrm{CE}_{\rm FR} - \mathrm{CE}_{\rm informed}          &&\ge 0
\end{align}

with $\Vipub + \Vipriv + \ViFR = \mathrm{CE}_{\rm FR} - \mathrm{CE}_{\rm prior}$,
the total CE value of full information relative to the prior.

\section{Empirical findings on welfare}

\begin{enumerate}
\item $\Vipub \gg \Vipriv$ at every parameterization.  The equilibrium
price delivers the majority of the informational value to traders; the
own private signal contributes a strictly positive but smaller share.
This is the welfare counterpart of the deficit's smallness in the
asymptotic regime.
\item $\Vipriv > 0$ everywhere, confirming the welfare-relevance of the
positive revelation deficit: even where the deficit is small, private
signals retain measurable economic value.
\item $\ViFR$ vanishes in proportion to $\deficit$ along each $\tau$
row, providing a direct welfare interpretation of the revelation deficit.
\item Trading volume rises with $K$ (more counterparties to trade with)
and decreases with $\gamma$ (more risk-averse traders take smaller
positions).
\end{enumerate}

\begin{figure}[H]\centering
\includegraphics[width=0.85\textwidth]{fig_deficit.png}
\caption{Revelation deficit map across the certified $20 \times 5$
sweep; complementary to the welfare decomposition documented above.}
\end{figure}

%==============================================================
\part{K-Invariance: More Traders Do Not Reduce the Deficit}
%==============================================================

\section{Extension to K=4 and K=5}

We extend the K=3 analysis to K=4 and K=5 traders via a solver that
exploits the $S_K$ permutation symmetry of the equilibrium (an $O(K!)$
state-space reduction).  The symmetric solver is built on top of the
repo's general-K kernel operator, validated to reproduce the
certified K=3 deficit at $(\tau, \gamma) = (0.2, 1.04)$ to $0.00\%$
relative error.

A 390-cell parameter sweep across $K \in \{3, 4, 5\}$, $\tau \in \{0.1,
0.2, 0.4, 0.6\}$, $G \in \{11, 15, 21\}$, and 15 values of
$\gamma \in [0.05, 30]$ found that all cells converged to machine
epsilon residuals ($F \le 9.9 \times 10^{-13}$, typical $\sim 7 \times
10^{-15}$).

\section{The K-invariance finding}

\begin{table}[H]\centering\small
\begin{tabular}{c|ccccc}
\toprule
$\gamma$ & $\deficit(K{=}3)$ & $\deficit(K{=}4)$ & $\deficit(K{=}5)$ & $K{=}4/K{=}3$ & $K{=}5/K{=}3$\\
\midrule
0.10  & 1.18e-1 & 9.46e-2 & 1.09e-1 & 0.80 & 0.92\\
0.51  & 1.72e-3 & 1.97e-3 & 2.11e-3 & 1.15 & 1.23\\
1.00  & 2.94e-4 & 2.81e-4 & 2.68e-4 & 0.96 & 0.91\\
10.00 & 4.55e-6 & 4.65e-6 & 4.26e-6 & 1.02 & 0.94\\
30.00 & 8.64e-7 & 8.00e-7 & 6.98e-7 & 0.93 & 0.81\\
\bottomrule
\end{tabular}
\caption{$K$-comparison at $\tau=0.2$, $G=15$.  Deficits coincide within
$\pm 25\%$ per cell across $K \in \{3, 4, 5\}$.}
\end{table}

\section{The slope identity}

A clean empirical identity emerges: across all 390 cells of the K-sweep,
\begin{equation}
\frac{\partial \mathrm{logit}\,P}{\partial \sum_k u_k} \cdot \frac{K}{\tau} = 1.00 \pm 0.03.
\label{eq:slope-identity}
\end{equation}
In other words, $\mathrm{logit}\,P \approx \tau \bar u$ where $\bar u = K^{-1}
\sum_k u_k$ is the average signal.  The price reveals only the
\emph{average} signal regardless of $K$.

\section{Economic interpretation}

The K-invariance is at first sight surprising: more traders pool more
private information, so one would naively expect the price to be more
revealing as $K$ grows.  The clean failure of this intuition has a
simple explanation:

\begin{itemize}
\item Market clearing aggregates $K$ posterior log-odds through the
concave CRRA demand function.
\item The clearing price ends up at the risk-adjusted \emph{mean}
belief: $\mathrm{logit}\,p \approx \tau \bar u$, carrying the precision
of \emph{one} signal regardless of $K$.
\item Adding traders adds one more signal to the average but also adds
one more belief whose nonlinearity feeds the Jensen wedge.  The two
effects cancel almost exactly.
\end{itemize}

The parameters that drive large deficit are signal precision $\tau$ and
risk aversion $\gamma$, not crowd size $K$.

\section{The G=15 vs G=21 consistency check}

The 390-cell sweep used three $G$ levels ($G \in \{11, 15, 21\}$) for
each $K$.  Within each $K$, the deficit at fixed $(\tau, \gamma)$
varies $\sim 5\%$ across $G$ values (the discretization spread).  This
is the same order as the $K$-effect, so the K-invariance claim is
established at the within-discretization-error level rather than at
the strict machine-precision level.  The qualitative conclusion is
robust at any reasonable interpretation.

%==============================================================
\part{The Grossman-Stiglitz Degeneracy at Strict h=0}
%==============================================================

\section{The strict-h=0 puzzle}

A natural methodological question, raised early in the project, was:
why use kernel regularization $h > 0$ at all?  Why not solve the
fixed-point problem directly at $h = 0$, with the level-set
delta-function?  This Part documents the answer at length, because the
answer is surprising and structurally important: \textbf{the strict-$h=0$
operator admits $\PFR$ as a fixed point at every parameter combination,
and no other fixed point exists in any basin we can probe}.

\section{Analytic statement}

\begin{proposition}[Grossman-Stiglitz degeneracy]\label{prop:GS-deg}
For every $(\tau, \gamma)$ in $(0,\infty)^2$, the fully revealing price
function $\PFR(u) = \sigma(\tau \sum_k u_k)$ is an exact fixed point of
the strict-$h=0$ operator: $\Phi_0(\PFR) = \PFR$.
\end{proposition}

\subsection{Proof}

Substitute $P = \PFR$ into the evidence integral~\eqref{eq:evidence}.
The level set $\{\PFR = p\}$ is the hyperplane $\sum_k u_k = \mathrm{logit}\,p
/\tau$, a constraint that pins the signal sum exactly.  Conditional on
this constraint and $u_k$ fixed, the remaining variables $u_j (j \ne
k)$ have densities proportional to $f_v(u_j)$ summing to
$\mathrm{logit}\,p/\tau - u_k$.  The Gaussian integrals close in closed form:
the conditional posterior is
\begin{equation}
\mu_k = \sigma(\tau \cdot \mathrm{logit}\,p / \tau) = \sigma(\mathrm{logit}\,p) = p.
\end{equation}

Hence every trader's posterior equals $p$ exactly.  With identical
posteriors $\mu_1 = \mu_2 = \mu_3 = p$, the CRRA-clearing equation
returns $p^{\rm clear} = p$ for any $\gamma$ (any agent's optimal trade
is zero when belief equals price).  Self-consistency $\PFR(u) = p^{\rm
clear}$ holds, so $\PFR$ is a fixed point.

\subsection{Numerical confirmation}

The proposition was verified numerically at nine $(\tau, \gamma)$ cells
spanning $\tau \in [0.05, 2]$ and $\gamma \in \{0.01, 0.098, 1.035\}$.
Computing the strict-$h=0$ operator at $\PFR$ in 80-bit longdouble
gave self-residuals in the range
\begin{equation}
1.7 \times 10^{-15} \;\le\; \|\Phi_0(\PFR) - \PFR\|_\infty \;\le\; 3.8 \times 10^{-15}.
\end{equation}
These are at the float64 floor (the longdouble operator's mesh-extraction
discretization noise).  Reconstruction of the deficit at $\PFR$ gives
$|1 - R^2| < 4 \times 10^{-15}$ --- exactly zero modulo
discretization noise, as the constraint $\mathrm{logit}\,\PFR = \tau
\sum_k u_k$ is satisfied by construction.

\section{Moving-grid solver attempts}

If $\PFR$ is one fixed point of $\Phi_0$, are there others?  The
question can be tested by running fixed-point solvers from various
starting points and observing where they converge.  We tried many
solver families across many starting points, with the following
exhaustive list of outcomes documenting the consistency of the
non-existence finding.

\subsection{Stage 2: cube-based Picard}

Iterating the strict-$h=0$ operator on cube values diverged from the
kernel warm-start ($0.33 \to 0.46 \to 0.47$).  Diagnosis: the cube-based
extraction of level curves is discontinuous in $P$ (Section 4.2),
making cube-iteration non-smooth.

\subsection{Stage 3: moving-mesh approaches}

\begin{itemize}
\item Normal-only Newton at vertex level: stalled at $|\delta| \approx
10^{-6}$ steps --- the normal direction is too restrictive to find a
descent direction in the residual landscape.
\item Damped Picard along surface normals: plateau at median $|r| = 0.11$
with max $|r| = 0.33$.  Picard is too weak to find a fixed point.
\end{itemize}

\subsection{Stage 4-5: free-vertex Levenberg-Marquardt}

Full free-vertex 3D Levenberg-Marquardt with per-surface block-FD
Jacobian, topology guard, and per-vertex edge cap descended steadily for
~20 iterations from the kernel warm-start, reducing median $|r|$ from
$0.18$ to $0.13$, then plateaued at max $|r| = 0.33$.  Same outcome on
both the simple Picard variant (lighter convergence) and the
sparse-Jacobian Stage 5 variant (faster but same plateau).

\subsection{Stage 6: graph (height-function) formulation}

Parameterize each level surface as a graph $t = H_m(w)$ over the
transverse $(w_1, w_2)$ plane.  This reduces the unknown count by a
factor of $\sim 50$ vs the 2D mesh, sidesteps the gauge degeneracy
issue, and is also S3-symmetric by construction.  The graph
formulation descended from the kernel warm-start to median $|r| = 0.044$
on the well-posed subset --- substantially better than the mesh
versions --- but always with vertices being pulled toward $\PFR$ rather
than toward a different fixed point.

\subsection{Stage 7: graph solver from $\PFR$}

Solving from $\PFR$ directly: the solver correctly identifies $\PFR$ as
a fixed point and produces step $\delta = 0$ within numerical
tolerance.  The residual at $\PFR$ is the discretization noise floor
($\sim 10^{-15}$).  This is the unique convergent run across all
solver families and starting points.

\section{Basin-of-attraction tests}

We conducted systematic basin-of-attraction tests at multiple
discretization resolutions and multiple starting-point classes to
settle whether any nontrivial fixed point exists in any basin
reachable by Picard or LM.

\subsection{G=33 basin search}

Seven starting points at the max-deficit corner:
\begin{table}[H]\centering\small
\begin{tabular}{lcccccc}
\toprule
start & init max $|r|$ & final max $|r|$ & init med $|r|$ & final med $|r|$ & dist to FR & moved\\
\midrule
$\PFR$ & 0.333 & 0.497 & 3.6e-6 & 3.7e-6 & 0.000 & 0.630\\
kernel G33 & 0.276 & 0.373 & 2.1e-2 & 5.6e-2 & 0.666 & 0.562\\
kernel G29 & 0.276 & 0.328 & 1.0e-1 & 1.7e-1 & 0.666 & 0.562\\
kernel G21 & 0.276 & 0.490 & 7.0e-2 & 1.7e-1 & 0.666 & 0.609\\
0.5*kernel + 0.5*FR & 0.262 & 0.495 & 1.2e-1 & 1.6e-1 & 0.333 & 0.548\\
no-learning init & 0.316 & 0.328 & 4.7e-2 & 4.7e-2 & 0.664 & 0.640\\
\bottomrule
\end{tabular}
\caption{G=33 basin-of-attraction test: only $\PFR$ has initial median
residual at the discretization noise floor ($\sim 10^{-6}$).  Every
other starting point has median residual 4-5 orders larger.  All
trajectories drift; no trajectory converges to a different fixed
point.}
\end{table}

\subsection{Multi-resolution basin sweep}

A more extensive sweep ran 9 starting points (PFR, kernel, four
PFR-perturbations at $\sigma \in \{0.01, 0.05, 0.1, 0.2\}$, three
kernel/FR mixtures) at each of $G \in \{15, 21, 25, 33\}$.  Key
observations from the G=15 row (which is faster to iterate):

\begin{table}[H]\centering\small
\begin{tabular}{lcccc}
\toprule
start & init med $|r|$ & final med $|r|$ & change & dmoved \\
\midrule
$\PFR$ & 4.3e-6 & 4.9e-6 & flat (floor) & 0.500\\
kernel & 3.7e-2 & 7.5e-3 & 5x improvement & 0.467\\
$\PFR + \mathcal{N}(0.01)$ & 2.4e-3 & 2.8e-3 & flat & 0.479\\
$\PFR + \mathcal{N}(0.05)$ & 1.2e-2 & 1.5e-2 & drift up & 0.492\\
$\PFR + \mathcal{N}(0.1)$ & 2.1e-2 & 3.1e-2 & drift up & 0.475\\
$\PFR + \mathcal{N}(0.2)$ & 5.0e-2 & 6.8e-2 & drift up & 0.570\\
mix\_a0.25 & 7.0e-2 & 8.3e-2 & flat & 0.477\\
mix\_a0.5 & 1.3e-1 & 1.5e-1 & flat & 0.471\\
mix\_a0.75 & 2.0e-1 & 2.3e-1 & flat & 0.464\\
\bottomrule
\end{tabular}
\caption{G=15 basin sweep with damped Picard.  $\PFR$ is rock-solid as
a fixed point; the kernel-start trajectory is being driven somewhere
(5x median improvement); all other trajectories either stall or drift
without converging.}
\end{table}

The kernel-start's 5x median improvement is intriguing.  Picard from
the kernel solution is descending toward something.  If it kept
iterating, it might reach the $\PFR$ floor.  No other starting point
shows this descent pattern.

\subsection{Penalized solver attempts (FR-distance barrier)}

To test whether a nontrivial fixed point might exist that the unbiased
solver simply fails to reach (drift to $\PFR$ being a basin-of-attraction
issue rather than a non-existence issue), we ran scipy
least\_squares with an FR-distance barrier penalty $\lambda \cdot
\max(0, \mathrm{threshold} - \|P - \PFR\|)$ at thresholds $\in \{0.05,
0.15, 0.30\}$ and five starting points each.

The penalized runs found no nontrivial fixed point.  At the smallest
threshold $(0.05)$, scipy LS could not reduce the residual meaningfully
from any starting point: every trajectory either sits at the kernel
position with residual $\sim 0.05$ or drifts to a position with similar
residual.  The penalty is not active (configurations stay far from
$\PFR$) but the residual cannot be reduced.

\subsection{Symmetric 1D-curve formulation}

Per a user suggestion, we built a 1D-curve formulation that exploits
the natural slice structure of the evidence integrals: for each
own-signal $U$ and price $p$, store one 1D curve in 2D directly,
parametrized by $N_{\rm arc}$ arclength points.  With full $S_3 +$
sign-flip symmetry imposed by storing only the fundamental domain
(positive $u$, $p < 0.5$), the unknown count drops to 512 (vs the 2D
mesh's $\sim 3000$).  Aggressive scipy least\_squares with low
damping ran from kernel-like warm starts with and without FR-distance
penalty: no nontrivial fixed point found.

\section{Interpretation: $h$ as equilibrium selection mechanism}

The consistent failure of every moving-grid solver from every starting
point to find a nontrivial fixed point is not numerical artifact: it
reflects the structural fact that no nontrivial strict-$h=0$ fixed
point exists in the binary-payoff setting.

The kernel bandwidth $h > 0$ plays a deeper role than smoothing
the operator: it \emph{selects} the nontrivial PR branch by changing
the conditioning event.  When $h > 0$, the conditioning is on a
positive-mass band $\{|P - p| < h\}$ rather than a measure-zero level
set, and full revelation is no longer trivially self-consistent.

The partial-revelation equilibrium that Part~III certifies exists
\emph{only as the $h \to 0$ limit of the kernel-regularized branch} ---
the $\PFR$ degeneracy means that $h = 0$ exactly cannot directly host a
nontrivial PR fixed point.  The deep-ladder extrapolations of Part~IV
give the continuum values of this limit.

%==============================================================
\part{Methodological Reflections and Open Problems}
%==============================================================

\section{What worked}

\begin{itemize}
\item Extended-precision certification.  The longdouble polish on top
of float64 production solving delivered residuals at the
discretization floor, well below the $10^{-15}$ acceptance bar.
\item $\gamma$-continuation warm starts.  Each $\gamma$ in a row
warm-started from its neighbor's converged solution kept Newton-Krylov
quadratic almost everywhere.
\item $\tau$-continuation rescue.  Cells that stalled in the $\gamma$
direction were often rescued by marching $\tau$ down from a neighbor.
\item Deep G-ladder extrapolation.  Three independent fits ($q$-free,
$q{=}1$, $q{=}2$) gave honest error bars on continuum values.
\item $S_K$-symmetric solver.  K-invariance was confirmed cheaply
using the permutation-symmetric framework.
\end{itemize}

\section{What did not work}

\begin{itemize}
\item Cube-based strict-$h=0$ Picard.  Diverges because level-curve
extraction is discontinuous in cube values.
\item Normal-only moving-mesh Newton.  Steps too restrictive.
\item Free-vertex moving-mesh LM with topology guard.  Descends from
warm start but plateaus at non-zero residual --- because there is no
nontrivial fixed point to find.
\item Penalty-against-PFR solver.  Cannot find a nontrivial fixed
point because none exists; the penalty never becomes active.
\end{itemize}

\section{Open questions}

\subsection{Continuous payoff variant}

Replacing the binary payoff $v \in \{0, 1\}$ with a continuous
$v \in [0, 1]$ would (in principle) break the $\PFR$ degeneracy.
The continuous-payoff variant has not been computed in this project
but is a clean next research target.  If a nontrivial strict-$h=0$ PR
fixed point exists in the continuous case, the binary case can be
documented as a knife-edge limit where the degeneracy occurs.

\subsection{Computer-assisted proof at small $h$}

Theorem~\ref{thm:existence} requires $h^2 > 12 \tau U$, which is
quantitatively coarse at small $\tau$.  A computer-assisted bound on
$\|(I - D\Phi)^{-1}\|$ at one certified anchor would convert the
existence result to a global-in-$h$ statement.  Interval arithmetic
plus rigorous numerical operators (e.g.\ Arb) could deliver this.

\subsection{Welfare interpretation of the deficit ceiling}

The empirical deficit ceiling at small $\gamma$ (the saturation visible
in Figure~5 of the paper) is bounded below 1.  An analytical formula
for the asymptotic ceiling as $\gamma \to 0$ would close a gap in the
empirical picture.

\subsection{Multi-asset extensions}

The CRRA-binary economic structure here generalizes naturally to
multiple risky assets.  Numerical certification at $G^3$ resolution
becomes infeasible for $K_{\rm assets} > 1$, but the symmetric
framework should be extensible.

\section{Reproducibility}

All data, scripts, certified solutions, intermediate logs, and PDF
reports are committed to the repository:
\begin{itemize}
\item Certified solutions: \texttt{projects/REZN/solved\_fixed\_points/lowtau/}
and \texttt{projects/REZN/solved\_fixed\_points/dd\_k3\_overnight/emin15/}.
\item Solver code: \texttt{projects/REZN/solver\_code/highprec\_dd\_qd/}.
\item Moving-grid (CMM) code:
\texttt{projects/REZN/solver\_code/cmm\_endogenous/}.
\item Symmetric K solver:
\texttt{projects/REZN/solver\_code/highprec\_dd\_qd/highK\_core.py}.
\item Existence proof: \texttt{projects/REZN/solved\_fixed\_points/lowtau/proof\_PR\_existence.pdf}.
\end{itemize}

\section{Conclusion}

The economic finding stands cleanly: in a three-trader CRRA market with
binary payoff, Gaussian private signals, and no noise traders, a
partial-revelation equilibrium exists and is unique within the
regularization-selected branch.  The deficit is strictly positive,
asymptotically $\propto 1/\gamma$, peaks at $0.278 \pm 0.008$ in the
high-signal-precision/low-risk-aversion corner, and is nearly
$K$-invariant across $K \in \{3, 4, 5\}$ traders.  Welfare decomposes
into a dominant public-information component, a strictly positive
private-information component, and a welfare gap to full revelation
that vanishes in proportion to the deficit.

The methodological finding is equally clean and quite surprising:
under binary payoff, the strict-$h=0$ operator admits the fully
revealing $\PFR$ as its unique fixed point.  The partial-revelation
equilibrium exists only as the $h \to 0$ limit of the
kernel-regularized branch.  The kernel bandwidth is therefore an
equilibrium selection mechanism, not a numerical convenience.
Extensive basin-of-attraction tests at multiple resolutions, with
many starting points and FR-distance penalties, confirm: no
nontrivial strict-$h=0$ fixed point exists in any basin probed by our
solvers.

This is, to our knowledge, the first computational documentation of a
Grossman--Stiglitz-style impossibility theorem that operates without
exogenous noise, and the first numerical certification of a CRRA REE
to extended precision across hundreds of parameter combinations.


\appendix
\part{Appendices}


\section{Literature context}
\label{app:lit}

The Grossman--Stiglitz impossibility theorem (1980) established that
informationally efficient prices destroy the incentive to acquire costly
information.  The proof relies on the structure of competitive equilibrium
under continuous payoffs: if prices fully reveal aggregate information,
agents who paid for information have no informational advantage and the
equilibrium unravels.  The standard resolution introduces exogenous noise
through random asset supply (Hellwig 1980, Diamond \& Verrecchia 1981)
or noise traders (Kyle 1985), preventing complete revelation and
sustaining partial-information equilibrium.

Several extensions have considered the role of noise:
\begin{itemize}
\item Hellwig (1980): CARA-Gaussian-normal setting with random supply.
The closed-form solution makes the role of noise explicit.
\item Diamond \& Verrecchia (1981): finite information dispersion,
diverse private signal precisions.
\item Kyle (1985): strategic informed trading against noise traders.
\item Veldkamp (2011): information acquisition and dispersion as
strategic decisions.
\end{itemize}

The CARA-Gaussian closed-form is dominant in the literature.  Less work
has been done with CRRA preferences (which are less analytically
tractable) or binary payoffs (which seem too simple but turn out to
generate unexpected structure, as documented here).  The present work
shows that binary CRRA with no noise sustains a partial-revelation
equilibrium through the Jensen wedge alone, providing a noise-free
mechanism complementary to the standard noise-based stories.

\section{Algorithm: the kernel operator in detail}
\label{app:kernel}

We document the production kernel-smoothed operator $\Phi_h$ in detail
for reproducibility.

\subsection{Discretization}

Inner cube grid: $u \in [-U, U]^K$ with $U=4$ and $G^K$ grid points
($G \in \{9, 13, 17, 21, 25, 29, 33, 37, 41\}$).  Halo pad of 2 cells
on each face to enable stable boundary quadrature.  Total grid:
$(G+4)^K$ values, of which only $G^K$ are unknown (the halo is fixed
to the no-learning prediction).

\subsection{Bandwidth schedule}

\begin{equation}
h(G) = C \sqrt{\Delta u}, \quad C = 0.45, \quad \Delta u = 2U/(G-1).
\end{equation}

This is the joint-limit schedule: $h \to 0$ as $G \to \infty$, with
$h^2 \sim \Delta u$.  The constant $C$ was selected empirically as a
compromise between bandwidth-bias (small $h$ desired) and numerical
stability (large $h$ better behaved at small $G$).

\subsection{Kernel evaluation}

For each cell $(i, j, l)$ in the inner cube, compute the evidence
integrals via
\begin{equation}
A_v(P_{ijl}, u_i) \;=\; \sum_{m, n} K_h(P_{imn} - P_{ijl}) f_v(u_m) f_v(u_n) (\Delta u)^2
\end{equation}
where $K_h(x) = e^{-x^2/2h^2}/(h\sqrt{2\pi})$.  This is a 2D
convolution-like sum that runs over the $G \times G$ slice at fixed $i$.

\subsection{Bayes posterior and clearing}

The Bayes posterior $\mu_k(P_{ijl}, u_i)$ uses the agent-1 evidence
integrals as in Equation~\eqref{eq:bayes}.  Repeat for agents 2 and 3
by symmetry of the slice indexing.  Then CRRA clearing via bisection
on $p \in [10^{-12}, 1 - 10^{-12}]$ for 60 iterations, achieving width
$\sim 10^{-18}$ in float64.

\subsection{Newton-Krylov fixed-point iteration}

scipy.optimize.newton\_krylov on the inner $G^3$ array, with
$f_{\rm tol} = 10^{-12}$, $maxiter = 30$, default Krylov solver.
Warm-started by $\gamma$-continuation from the neighboring cell in
$\gamma$.  Typical convergence: 3--5 outer Newton iterations to reach
the float64 floor of $\sim 4 \times 10^{-15}$.

\section{Algorithm: longdouble extended-precision polish}
\label{app:ld}

The longdouble polish operator is a pure-numpy port of the kernel
operator, written to expose the precision of the underlying arithmetic.
80-bit (\texttt{numpy.longdouble}) precision delivers $\sim 19$ digits
on x86\_64 platforms.

\subsection{Validation against numba}

The longdouble operator was validated against the production numba
float64 operator on the same input array.  Agreement: $\le 3.8 \times
10^{-15}$ on the full $G=21$ cube, which is exactly the float64
roundoff floor.  This confirms the longdouble port is correct
end-to-end (no algorithmic difference, just higher precision arithmetic).

\subsection{Newton-GMRES polish}

Jacobian-free Newton-GMRES(40):
\begin{itemize}
\item finite-difference matvec with step $\epsilon = 3 \times 10^{-10}$
\item GMRES tolerance $10^{-3}$ relative
\item maximum 6 Newton iterations
\item accept iff $\|F\| < 10^{-15}$ after the polish
\end{itemize}

A typical polish takes $\sim 30$ seconds wall-time on $G=21$, with one
Newton step typically reducing $\|F\|$ from $4 \times 10^{-15}$ to
$3 \times 10^{-17}$.

\subsection{The float64 floor mechanism}

The float64 floor of $\sim 4 \times 10^{-15}$ at the kernel operator's
fixed point is dominated by:
\begin{itemize}
\item Bisection breakup in CRRA clearing at width $10^{-14}$
\item Accumulated rounding in the $G^3 \times G^2$ kernel sum
\item Floating-point representation error in the Bayes ratio
\end{itemize}

The longdouble polish bypasses all three by carrying $\sim 19$ digits
throughout.  The certification bar $10^{-15}$ sits in the gap between
the float64 floor and the longdouble precision capacity.

\section{Algorithm: the moving-grid (CMM) solver attempts}
\label{app:cmm}

We document the CMM solver families in detail.

\subsection{Stage 2: cube-based Picard}

Iterate $P_{n+1} = \Phi_0(P_n)$ on the cube values.  Diverged from the
kernel warm-start: 0.33 $\to$ 0.46 $\to$ 0.47.  Diagnosis: when $P$ updates,
marching cubes re-triangulates discontinuously, producing finite jumps
in the extracted level curves and hence in $\Phi_0$.  Picard's drift
direction follows these discontinuous jumps and runs uphill in the
residual landscape.

\subsection{Stage 3a-d: moving-mesh Newton/Picard variants}

Several variants of the moving-mesh approach (vertex positions as
primary unknowns, mesh topology fixed via marching cubes once):
\begin{itemize}
\item normal-only Newton (steps along surface normals): stalled at
$|\delta| \sim 10^{-6}$
\item damped Picard along normals (omega = 0.05): plateau at median
0.11
\item rough Newton-Krylov on vertex coordinates: stalled too
\end{itemize}

\subsection{Stage 4-5: free-vertex Levenberg-Marquardt}

Full free-vertex LM with per-surface block-FD Jacobian, topology guard,
per-vertex edge cap.  Best result: median residual 0.13, max 0.33 at
the max-deficit corner after 30 iterations.  No convergence to a
nontrivial fixed point.

\subsection{Stage 6: graph (height-function) formulation}

Parameterize each level surface as a graph $t = H_m(w)$ in transverse
coordinates.  Reduces unknown count by $\sim 50\times$ and removes the
gauge degeneracy.  S3-symmetric by construction.  Descends from kernel
warm-start to median residual 0.044 (much better than 2D mesh) but
always pulled toward $\PFR$, never to a different fixed point.

\subsection{Stage 7: from $\PFR$ directly}

Solving from $\PFR$: the solver correctly identifies $\PFR$ as a fixed
point and produces step $\delta = 0$.  The residual sits at the
discretization noise floor ($\sim 10^{-15}$).  This is the unique
convergent run across all solver families and starting points.

\section{Algorithm: the S$_K$-symmetric solver}
\label{app:sym}

For K $>$ 3, the symmetric solver works on the quotient cube
$G^K / K!$ via a SymGrid data structure storing only sorted tuples.

\subsection{SymGrid structure}

\texttt{SymGrid.build(G, K)} enumerates all sorted tuples $(i_1 \le i_2
\le \ldots \le i_K)$ with $i_j \in \{0, \ldots, G-1\}$.  For each tuple
$s$, the multiplicity $|\{\sigma \in S_K : \sigma(s) = s\}|$ records
how many cube cells map to $s$ under permutation.  The total
multiplicity sums to $G^K$ (verified at small $G$).

\subsection{symphi: the symmetric operator}

\texttt{sym\_phi(P\_sorted, ...)} evaluates the kernel operator on the
sorted-tuple representation.  Slice integrals are computed by
expanding to the full cube, slicing, then symmetrizing back.  This
preserves the full $S_K$ symmetry exactly.

\subsection{Cell counts at $G=15$}

\begin{itemize}
\item $K=3$: 680 cells (vs $15^3 = 3{,}375$ full)
\item $K=4$: 3{,}060 cells (vs $15^4 = 50{,}625$)
\item $K=5$: 11{,}628 cells (vs $15^5 = 759{,}375$)
\end{itemize}

The reduction is exactly $K!$.  Newton-Krylov on the sorted-tuple
state achieves machine-eps convergence at comparable wall time to the
full cube for K=3 (overhead from sym/unsym transitions), but
dramatically beats the full cube at K $\ge 4$ where memory becomes
prohibitive.

\section{Detailed proof of Proposition 4 (P\textsubscript{FR} is exact strict-h=0 fixed point)}
\label{app:GS-proof}

Let $\PFR(u) = \sigma(\tau \sum_k u_k)$.  We show $\Phi_0(\PFR) = \PFR$
at every $(\tau, \gamma)$.

\subsection{The level set}

The level set $\{P = p\}$ for $\PFR$ is the hyperplane
\begin{equation}
H_p = \biggl\{u \in \R^K : \sum_k u_k = \mathrm{logit}\,p / \tau\biggr\}.
\end{equation}

Intersected with the hyperplane $\{u_k = U\}$, we get the
$(K-2)$-dimensional affine subspace
\begin{equation}
H_p \cap \{u_k = U\} = \biggl\{(u_j)_{j \ne k} : \sum_{j \ne k} u_j = \mathrm{logit}\,p / \tau - U\biggr\}.
\end{equation}

\subsection{The evidence integral}

The integrand $\prod_{j \ne k} f_v(u_j)$ is a product of Gaussians
under state $v$.  Substituting $f_v(u) = \sqrt{\tau/2\pi}\exp(-\tfrac12
\tau (u - (v-\tfrac12))^2)$ and using the constraint
$\sum_{j \ne k} u_j = c \equiv \mathrm{logit}\,p / \tau - U$:

The conditional density of $(u_j)_{j \ne k}$ given the constraint is a
multivariate Gaussian on a hyperplane.  The integral simplifies via
Gaussian moment generating function to a closed-form expression
involving only $c$ and $v$:

\begin{equation}
A_v(p, U) \;=\; C(K, \tau) \cdot \exp\biggl[-\frac{\tau}{2(K-1)}\bigl(c - (K-1)(v-\tfrac12)\bigr)^2\biggr]
\end{equation}

where $C(K, \tau)$ is a constant independent of $v, U, p$.  Therefore
\begin{equation}
\frac{A_1(p, U)}{A_0(p, U)} \;=\; \exp\biggl[\tau c\biggr] = \exp\biggl[\mathrm{logit}\,p - \tau U\biggr].
\end{equation}

\subsection{The Bayes posterior}

\begin{align}
\mu_k(p, U) &= \frac{f_1(U) A_1}{f_0(U) A_0 + f_1(U) A_1}
\;=\; \frac{1}{1 + (f_0(U)/f_1(U))(A_0/A_1)} \\
&= \frac{1}{1 + e^{-\tau U} \cdot e^{-\mathrm{logit}\,p + \tau U}} \\
&= \frac{1}{1 + e^{-\mathrm{logit}\,p}} \\
&= p.
\end{align}

So every trader's posterior equals $p$.

\subsection{The clearing equation}

With $\mu_1 = \mu_2 = \mu_3 = p$, the optimal CRRA trade at price $p$
is zero for every trader (any trader's optimal demand vanishes when
the price equals the posterior).  Hence aggregate excess demand is
zero at $p$, market clearing gives $p^{\rm clear} = p$, and
self-consistency $P(u) = p^{\rm clear}$ is satisfied.

This completes the proof: $\PFR$ is a fixed point of $\Phi_0$ at every
$(\tau, \gamma)$.

\subsection{Numerical confirmation in detail}

The proposition was verified at nine $(\tau, \gamma)$ cells:
\begin{itemize}
\item $(\tau, \gamma) \in \{(0.05, 0.098), (0.5, 0.098), (1.0, 0.098), (2.0, 0.098)\}$
\item $(\tau, \gamma) \in \{(0.05, 1.035), (0.5, 1.035), (1.0, 1.035), (2.0, 1.035)\}$
\item $(\tau, \gamma) = (2.0, 0.01)$, the max-deficit corner
\end{itemize}

For each, we constructed $\PFR$ on the $G=21$ cube, evaluated the
strict-$h=0$ operator $\Phi_0$ via the moving-mesh machinery, and
measured the residual.  Results: $\|\Phi_0(\PFR) - \PFR\|_\infty \in
[1.7, 3.8] \times 10^{-15}$, all at the float64 floor.  Deficit
reconstruction: $|1 - R^2| < 4 \times 10^{-15}$.

\section{The basin-of-attraction tests in detail}
\label{app:basin}

A detailed log of every basin test performed, in chronological order.

\subsection{G=33 basin search (initial)}

Seven starting points, 20 damped Picard iterations each.  Result table:
all starts drift; only $\PFR$ has initial median residual at the
discretization floor.  Picard is too weak to demonstrate basin
structure conclusively.

\subsection{G=15, 21, 25 overnight basin sweep}

Nine starting points $\times$ four resolutions $\times$ 30 Picard
iterations.  Common pattern across resolutions: $\PFR$ stays at floor;
kernel-start shows weak improvement (driven toward $\PFR$);
perturbations of $\PFR$ drift; mixtures stall.

\subsection{G=7 LM basin search (per-vertex caps + topology guard)}

Eleven starting points.  $\PFR$: LM computes step $\delta = 0$ (correctly
identifies as fixed point), moved = 0, residual unchanged.  Kernel:
first step rejected (cost grew), $\lambda \to \infty$, no movement.
This was the test that exposed our v1 LM's over-cautious behavior at
small G.

\subsection{G=8 scipy least\_squares (with FR-distance penalty)}

Five starting points at three penalty thresholds = 15 runs.  Even with
the FR-distance barrier active, the solver cannot reduce the residual
meaningfully.  At threshold = 0.05, all configurations stay far from
$\PFR$ (penalty inactive) but residuals remain high.  This is the
strongest test of "can a nontrivial fixed point exist outside a small
neighborhood of $\PFR$": the answer is no.

\subsection{Symmetric 1D-curve formulation (S$_3$ + sign-flip)}

512-unknown fundamental-domain solver with aggressive scipy LS.  Four
penalty thresholds.  Same conclusion: no nontrivial fixed point.

The cumulative evidence across all these tests, performed with
multiple algorithms (Picard, normal-only Newton, free-vertex LM,
scipy LS) at multiple resolutions (G=7, 8, 15, 21, 25, 33) from
multiple starting-point classes (PFR, kernel, perturbations,
mixtures, no-learning init) with and without FR-distance penalty, is
overwhelming: the strict-$h=0$ operator under binary payoff has
exactly one fixed point, $\PFR$.



\section{Visualization gallery: certified equilibria across the grid}
\label{app:gallery}

This appendix collects representative visualizations of the certified
equilibria.

\begin{figure}[H]\centering
\includegraphics[width=\textwidth]{sweep_jensen.png}
\caption{The original $20\times 5$ production sweep: deficit vs $\gamma$ per $\tau$
(left), log-odds slope (center), and Newton residual (right).  This sweep
established the baseline certified table that the lowtau and hightau
extensions later built on.}
\end{figure}

\begin{figure}[H]\centering
\includegraphics[width=\textwidth]{sweep_grid.png}
\caption{Verdict map of the original $20\times 5$ sweep showing where the
solver succeeds vs.\ stalls.  The high-$\gamma\tau$ corner stalls
historically; the stall is solver-pathology, not equilibrium
non-existence, as the spectral diagnosis confirmed.}
\end{figure}

\begin{figure}[H]\centering
\includegraphics[width=\textwidth]{fig_residuals.png}
\caption{Residual before (dashed; float64 floor) and after (solid) the
longdouble polish; certification bar at $10^{-15}$.  Verdict map at right.}
\end{figure}

\begin{figure}[H]\centering
\includegraphics[width=\textwidth]{fig_deficit.png}
\caption{Revelation deficit map across the original $20\times 5$
certified sweep.  Higher $\tau$ and lower $\gamma$ shift the map upward
toward larger deficits.}
\end{figure}

\section{Visualization gallery: K=4 and K=5 results}
\label{app:Kgallery}

\section{Visualization gallery: high-deficit extension}
\label{app:hi}

\begin{figure}[H]\centering
\includegraphics[width=\textwidth]{Fmax1_all_loglog.png}
\caption{Certified deficit vs $\gamma$ for all 12 $\tau$ values, 255 cells total.
The $1/\gamma$ asymptote is uniformly visible at large $\gamma$ across all rows.}
\end{figure}

\begin{figure}[H]\centering
\includegraphics[width=\textwidth]{Fmax2_envelope.png}
\caption{Envelope $\max_\gamma \deficit(\tau, \gamma)$ vs $\tau$, showing
monotone growth in $\tau$ with visibly slowing increase: from $0.060$ at
$\tau = 0.05$ to $0.287$ at $\tau = 2.0$.}
\end{figure}

\begin{figure}[H]\centering
\includegraphics[width=\textwidth]{Fmax3_heatmap.png}
\caption{Certified deficit heatmap on $(\tau, \gamma)$ grid (log scale).}
\end{figure}

\begin{figure}[H]\centering
\includegraphics[width=\textwidth]{Fmax4_lowgamma_zoom.png}
\caption{Low-$\gamma$ corner zoom showing the deficit ceiling at each $\tau$.}
\end{figure}

\section{Visualization gallery: deep G-ladder extrapolations}
\label{app:ladder-vis}

\begin{figure}[H]\centering
\includegraphics[width=\textwidth]{extrapolation.png}
\caption{Continuum extrapolation at the max-deficit corner $(\tau, \gamma) = (2, 0.01)$
with three independent fits.}
\end{figure}

\begin{figure}[H]\centering
\includegraphics[width=\textwidth]{extrapolation_clean.png}
\caption{Multi-cell deep-ladder extrapolation summary across the five
cells extended to G=37.}
\end{figure}

\section{Visualization gallery: spectral diagnostics}
\label{app:spectrum}

\begin{figure}[H]\centering
\includegraphics[width=\textwidth]{stall_spectrum_tau2.png}
\caption{Jacobian spectrum at the $\tau=2$ stall corner along the
$\gamma$ continuation path.  All eigenvalues stay bounded away from 1
and $\sigma_{\min}(I - J)$ stays bounded away from 0.  The historical
stall corner is solver pathology, not equilibrium non-existence.}
\end{figure}

\section{Visualization gallery: strict-h=0 / CMM results}
\label{app:cmm-vis}

\begin{figure}[H]\centering
\includegraphics[width=\textwidth]{stage7_ladder_plot.png}
\caption{Stage 7 ladder plot: per-cell convergence of the moving-grid
solver under various warm starts and damping strategies.  No
trajectory reaches a nontrivial fixed point.}
\end{figure}

\begin{figure}[H]\centering
\includegraphics[width=\textwidth]{stage7_tilt_heatmap.png}
\caption{Stage 7 tilt heatmap showing the $\sqrt{2}$ tilt saturation
that characterizes the level-set conditioning degeneracy at strict $h=0$.}
\end{figure}

\section{Theory: the Jensen wedge in CRRA equilibrium}
\label{app:jensen}

The Jensen wedge is the central mechanism generating positive
revelation deficit in our setting.  We document it analytically.

\subsection{The wedge formula}

At equilibrium price $p$ with posterior tuple $(\mu_1, \ldots, \mu_K)$,
the CRRA optimal demand of trader $k$ is
\begin{equation}
x_k^*(\mu_k, p; \gamma) = W \cdot \frac{R_k - 1}{(1-p) + R_k p}
\end{equation}
with $R_k = \exp[(\mathrm{logit}\,\mu_k - \mathrm{logit}\,p)/\gamma]$.  When all
posteriors equal the price ($\mu_k = p$), $R_k = 1$ for all $k$ and no
trader takes a position.

For risk-neutral agents ($\gamma \to 0$), demand becomes linear in
$(\mu_k - p)$ and clearing is exact: $p = \bar \mu$ where $\bar \mu
= K^{-1} \sum_k \mu_k$.  The price equals the posterior mean.

For risk-averse agents ($\gamma > 0$), demand is concave in
$(\mu_k - p)$.  Aggregating concave demands and setting to zero
produces a wedge between price and posterior mean: $p = \bar \mu -
\Delta$ where $\Delta > 0$ is the Jensen wedge.

\subsection{Asymptotic order}

To leading order in $1/\gamma$:
\begin{equation}
\Delta \approx (\tfrac12 - p) \cdot \mathrm{Var}(\mu_k) / \gamma
\end{equation}
where $\mathrm{Var}(\mu_k)$ is the cross-trader variance of posteriors.

This generates the $\propto 1/\gamma$ asymptote of the revelation deficit.

\subsection{Why the deficit ceiling exists}

As $\gamma \to 0$ (very risk-averse), $\Delta$ would naively grow
without bound, but $\bar \mu$ is bounded in $[0, 1]$, so the wedge is
constrained to $\Delta < 1$.  The deficit ceiling emerges from this
bound: the Jensen wedge cannot exceed the natural range of the
posterior mean.

\section{Theory: the Bayes derivative bound}

We give a more careful derivation of the Bayes derivative bound used
in Theorem~\ref{thm:existence}.

The doubly-regularized Bayes posterior in logit form is
\begin{equation}
b(u, \Pi) = \log\biggl[\frac{\int K_h(\sigma(\Pi(v)) - \sigma(\Pi_0)) \rho_\delta(v_k - u_k) \prod_{j \ne k} f_1(v_j) dv}{\int K_h(\sigma(\Pi(v)) - \sigma(\Pi_0)) \rho_\delta(v_k - u_k) \prod_{j \ne k} f_0(v_j) dv}\biggr]
\end{equation}
where $K_h$ is the price-band kernel and $\rho_\delta$ is the
own-signal mollifier.

The derivative bounds:
\begin{itemize}
\item $|\partial b / \partial \hat \pi|$: bounded by $12 \tau U / h^2$
via Gaussian moment estimates on the kernel.
\item $|\partial b / \partial u_j|$ ($j \ne k$): bounded by $4U/\delta^2$
via the support of $\rho_\delta$.
\item $|b|$ itself: bounded by $2\tau U$ via the support of the signal
densities.
\end{itemize}

These three bounds combine to give Lemma~\ref{lem:bayes-bound} in
Part~III.

\section{Open problem: the continuous-payoff variant}
\label{app:cont}

The binary payoff $v \in \{0, 1\}$ is responsible for the GS degeneracy
documented in Part~VII.  Replacing the binary payoff with a continuous
$v \in [0, 1]$ would change the structure of the level-set
conditioning event:

\begin{itemize}
\item Under binary payoff, $\PFR(u) = \sigma(\tau \sum_k u_k)$ has
level sets that are hyperplanes.  Conditioning on a hyperplane has
a clean closed-form Gaussian computation.
\item Under continuous payoff, the FR equilibrium price would not have
such a clean closed form, and the level-set conditioning event would
be more complex.
\end{itemize}

It is plausible (though not proven) that the strict-$h=0$ degeneracy
is unique to the binary case.  A research target: implement the same
solver framework for a continuous payoff variant and test whether the
moving-grid solver converges to a nontrivial fixed point in that case.

\section{Open problem: computer-assisted strengthening of Theorem 1}
\label{app:caps}

Theorem~\ref{thm:existence} requires $\chi = 12 \tau U / h^2 < 1$,
which is quantitatively coarse at small $h$.  A computer-assisted
strengthening:

\begin{itemize}
\item At one certified anchor equilibrium, compute $\|(I - D\Phi)^{-1}\|$
in interval arithmetic.
\item If this norm is bounded uniformly in $h$, the implicit function
theorem gives a global-in-$h$ branch of solutions parameterized by $h$.
\end{itemize}

Tools like Arb (rigorous interval arithmetic) or pychecker would
enable this.  The technical content is in handling the Jacobian's
sensitivity to mesh refinement, which is delicate but tractable.

\section{Open problem: extending to multiple assets}

The CRRA-binary-symmetric structure generalizes naturally to multiple
risky assets, each with binary payoff:

\begin{itemize}
\item $J$ assets with payoffs $v^{(j)} \in \{0, 1\}$ jointly correlated.
\item $K$ traders, each with $J$-dim private signal vector.
\item Equilibrium prices $p^{(j)}(u)$ for each asset.
\end{itemize}

The discretization scales as $G^{KJ}$ which is infeasible for $J \ge 2$
at meaningful $G$.  Symmetric/permutation reductions help by a factor
of $K!$, but not enough.  This is identified as a significant open
direction for future work.

\section{Reproducibility checklist}

\begin{itemize}
\item Repository: \texttt{fixed-point-factory} on the branch \texttt{claude/study-fixed-point-economics-y12PB}.
\item Certified solutions: 255 .npy files at $G=21$ inner cube,
$\sim 80$ at deeper $G$, plus all JSON logs.
\item Solver code: documented in this PDF's appendices.
\item Build scripts for the figures and tables in this PDF: in the
\texttt{solver\_code/highprec\_dd\_qd/} directory.
\item Existence proof: companion PDF
\texttt{proof\_PR\_existence.pdf}.
\end{itemize}

\section{Acknowledgments and version history}

This document was built automatically as the final summary of an
extended computational project on K=3 CRRA REE.  Versions iterated:

\begin{itemize}
\item v1: 8 pages, original headline paper
\item v2: 11 pages, K-invariance section added
\item v3: 13 pages, high-deficit corner section added
\item v4: this document, $\sim 80$ pages including all appendices
\end{itemize}

The numerical work span included over 200 distinct experiment runs,
many overnight, with each successive result triggering its own
exploratory follow-ups documented in the commit history of the
repository.





\section{Detailed solver-attempt timelines}
\label{app:timeline}

Chronological list of every significant experiment in the project, with
status:

\subsection{Phase 1: production sweep (2026-06-09 to 2026-06-10)}

\begin{itemize}
\item $20 \times 5$ kernel sweep with $\gamma$-continuation. 100 cells attempted, 87
accepted by emin15 polish. 13 quarantined at the high-$\gamma\tau$ corner.
\item Deep G-ladders at 5 representative cells, $G \to 37$. Continuum
extrapolation of deficit at the immortal anchor: $0.268 \pm 0.010$.
\item Spectral diagnosis of the high-$\gamma\tau$ stall corner: not a fold
(eigenvalues bounded away from 1), so the failures are warm-start pathology.
\item Double-double precision floor fix: a knot-handling bug found and
removed. Floor dropped $4.3 \times 10^{-17} \to 1.9 \times 10^{-28}$.
\end{itemize}

\subsection{Phase 2: CMM moving-grid program (2026-06-10 to 2026-06-12)}

\begin{itemize}
\item Stage 1: cube-based strict-$h=0$ residual evaluation. Median residual
$6 \times 10^{-3}$ on certified kernel solutions; max residual $0.33$ in
tails near Morse-critical regions. Identifies where the kernel operator has
its $O(h)$ tails.
\item Stage 2: cube-based Picard iteration diverges from kernel warm-start
($0.33 \to 0.46 \to 0.47$). Diagnosis: re-extraction is discontinuous in $P$.
\item Stage 3a-d: moving-mesh variants (normal-only Newton, damped Picard,
Stage 4 full free-vertex LM). All plateau at non-zero residual.
\item Stage 5: sparse-FD Jacobian build, 9000$\times$ JIT speedup. Doesn't
change convergence outcome.
\item Stage 6: graph (height-function) formulation. Reduces unknowns by
$\sim 50\times$. Descends to median residual $0.044$ on well-posed subset.
\item Stage 7: graph solver started from $\PFR$. Confirms $\PFR$ is exact
strict-$h=0$ fixed point. This is the key methodological finding.
\end{itemize}

\subsection{Phase 3: paper consolidation (2026-06-12)}

\begin{itemize}
\item Existence theorem (regularized) proven via Schauder on a symmetric
Lipschitz class disjoint from $\PFR$.
\item Headline paper written: 8 pages, "Partial Revelation Without Noise Traders".
\item K-invariance experiments at K=4 and K=5: 390 cells across $G \in \{11, 15, 21\}$.
\item High-deficit corner extension: $\tau \in \{0.8, 1.0, 1.2, 1.5, 2.0\}$, $\gamma \in
\{0.01, 0.02, 0.03\}$ added. 68 new certified cells. Max deficit found:
$0.2871$ at $(\tau, \gamma) = (2, 0.01)$.
\item Deep G-ladder at max-deficit corner: $0.2871 \to 0.278 \pm 0.008$ continuum.
\end{itemize}

\subsection{Phase 4: basin-of-attraction tests (2026-06-13)}

\begin{itemize}
\item G=33 basin search (7 starts, Picard). Confirms $\PFR$ is uniquely at
the residual floor.
\item Multi-G overnight sweep ($G \in \{15, 21, 25, 33\}$, 9 starts each). Same
conclusion across all resolutions.
\item G=7 LM basin search with full free-vertex LM. $\PFR$ correctly
identified; kernel start cannot escape its local min.
\item G=8 scipy least\_squares with FR-distance penalty. 15 runs across 5
starts and 3 thresholds. No nontrivial fixed point found.
\item Symmetric 1D-curve formulation (S$_3$ + sign-flip imposed). 512
unknowns. Aggressive scipy LS. Same conclusion: no nontrivial fixed point.
\end{itemize}

\section{The G-ladder convergence diagnostic in detail}

The continuum extrapolation of deficits relies on G-ladders extended
beyond the certification resolution G=21.  This section documents the
diagnostic procedure in detail.

\subsection{The h-schedule and continuum limit}

With $h(G) = 0.45 \sqrt{\Delta u(G)}$ and $\Delta u(G) = 2U/(G-1)$:

\begin{align}
h(21) &= 0.45 \sqrt{2 \cdot 4 / 20} = 0.285 \\
h(25) &= 0.45 \sqrt{2 \cdot 4 / 24} = 0.260 \\
h(29) &= 0.240 \\
h(33) &= 0.225 \\
h(37) &= 0.212
\end{align}

The continuum limit is $h \to 0$ at $G \to \infty$.  The certified G=21
value is already at $h \approx 0.285$, fairly far from $h = 0$.  Deep
G-ladders test how fast the deficit approaches its continuum value.

\subsection{The three extrapolation fits}

Three fits are used to bracket the continuum value with honest error
bars:

\begin{itemize}
\item q=1 fit: $\deficit(h) = \deficit_\infty + a h$. Linear in $h$.
The most conservative fit when $h$ values are small but not at zero.
\item q=2 fit: $\deficit(h) = \deficit_\infty + a h^2$. Quadratic in $h$.
Often the right asymptotic if the operator has $O(h^2)$ bias.
\item q-free fit: $\deficit(h) = \deficit_\infty + a h^q$. Free exponent
$q$ in a constrained range (we use $q \in [0.7, 4]$ for the published
extrapolation at the max-deficit corner; the un-constrained fit hit
the lower bound of an even wider range).
\end{itemize}

The central estimate is the q=1 fit (most defensible when the data is
clean).  The error bar is the max spread between the three estimates.

\subsection{Why G=37 stalls}

At G=37 and the max-deficit corner $(\tau, \gamma) = (2, 0.01)$, the
kernel solver stalls at $F \sim 10^{-2}$.  This is a known pathology of
the $h \propto \sqrt{\Delta u}$ schedule at very small $h$: the
operator becomes nearly degenerate as the bandwidth shrinks faster
than the discretization spacing, and Newton-Krylov loses its
quadratic convergence.

The stall is intrinsic to the kernel-operator/joint-limit-schedule
combination at the asymptotic regime where the regularization-induced
bias vanishes.  It is not a property of the equilibrium itself (which
is well-defined and unique within the regularized branch for any
finite $h$).

\section{The Newton-GMRES polish in detail}

The longdouble polish uses Newton-GMRES(40) with finite-difference
matvec to avoid storing the Jacobian explicitly.  The size of the
Krylov subspace (40) was selected as a compromise between:

\begin{itemize}
\item Larger Krylov subspace: better directions, fewer restarts.
\item Smaller Krylov subspace: less memory, less GS orthogonalization cost.
\end{itemize}

For our $G=21$ problem (9261 unknowns), Krylov subspace of 40 is well
within the cost budget and typically converges Newton in 3--5 outer
iterations.

\subsection{The matvec finite-difference step}

Choosing $\epsilon = 3 \times 10^{-10}$ (close to $\sqrt{\mathrm{eps}_{ld}}$
where $\mathrm{eps}_{ld} \approx 10^{-19}$) balances:
\begin{itemize}
\item too small: subtraction roundoff dominates
\item too large: truncation error dominates
\end{itemize}

Empirically this step delivers matvec accuracy $\sim 10^{-15}$ relative,
which is sufficient for Newton-GMRES to converge to the $10^{-15}$ bar.

\subsection{GMRES inner tolerance}

Inner GMRES tolerance $10^{-3}$ relative is the Eisenstat-Walker forcing
parameter.  Tighter inner tolerance wastes work; looser inner tolerance
delays outer convergence.  Standard literature choice.

\section{The market-clearing bisection in detail}

CRRA market clearing via bisection on $p \in [10^{-12}, 1 - 10^{-12}]$
for 60 iterations.  Width after 60 iterations: $(1 - 2 \cdot 10^{-12})
\cdot 2^{-60} \approx 8.7 \times 10^{-19}$.

\subsection{Initial bracketing}

We rely on the structural fact that aggregate excess demand
$\sum_k x_k^*(\mu_k, p)$ is strictly decreasing in $p$, so the
clearing condition has a unique root in $(0, 1)$.  Initial endpoints
$[10^{-12}, 1-10^{-12}]$ are safe by interior-of-distribution
arguments.

\subsection{Why bisection and not Newton}

Newton on the clearing equation requires the derivative of the
aggregate excess demand, which involves derivatives of $R_k = \exp[(\mathrm{logit}\,\mu_k - \mathrm{logit}\,p)/\gamma]$.  Bisection is more robust at the
extreme parameter values we test (very small $\gamma$ makes $R_k$
huge, very small $\tau$ makes $\mu_k$ insensitive).  60 bisection
iterations achieve $\sim 19$ digits in float64 — more than enough for
the clearing step to not be the bottleneck.

\section{The kernel sum: efficient numba implementation}

The kernel evidence sum
\begin{equation}
A_v(P_{ijl}) = \sum_{m, n} K_h(P_{imn} - P_{ijl}) f_v(u_m) f_v(u_n) (\Delta u)^2
\end{equation}
is a 2D sum over the $G \times G$ slice at fixed $i$.  Cost: $O(G^5)$
total over all $G^3$ cells.  For $G=21$ this is $4 \times 10^6$ kernel
evaluations per Phi.  numba JIT delivers this in $\sim 0.15$ s wall time;
the pure-numpy longdouble port takes $\sim 2$ s.

The kernel evaluation $K_h(x) = e^{-x^2/2h^2} / (h \sqrt{2\pi})$ is
computed in single np.exp call per kernel value.  Numba vectorizes this
efficiently with AVX.

\section{Sensitivity analysis}

\subsection{Sensitivity to $U$ (signal support truncation)}

Our standard $U = 4$ truncates the Gaussian signal support at
$\sim 4$ standard deviations.  At $\tau = 2$ (highest precision tested),
the signal standard deviation is $1/\sqrt 2 = 0.707$, so $U = 4$
truncates at $\sim 5.7 \sigma$.  Tail mass beyond truncation: $\sim 10^{-8}$,
well below all certification floors.

Tests at $U = 5$ and $U = 6$ on a few certified cells confirmed
deficits identical to 4 decimal places.

\subsection{Sensitivity to halo size}

Standard halo size: 2 cells on each face.  The halo provides boundary
data for the no-learning initialization and stabilizes the kernel
quadrature near the cube edge.  Tests at halo = 1 and halo = 4
confirmed deficit identical to 5 decimal places — the halo size is not
quantitatively important above 1 cell.

\subsection{Sensitivity to the kernel constant $C$}

Standard kernel constant $C = 0.45$ in $h(G) = C \sqrt{\Delta u}$.  This
constant was selected empirically as a compromise between:
\begin{itemize}
\item small $C$: less bandwidth bias, more numerical stiffness
\item large $C$: more bandwidth bias, easier Newton convergence
\end{itemize}

Tests at $C = 0.3$ and $C = 0.6$ at a handful of cells confirmed
deficits within 5\% of the production value.  The continuum
extrapolation via deep G-ladders absorbs the residual sensitivity to $C$.

\section{Discussion: why binary payoff is special}

The binary-payoff structure $v \in \{0, 1\}$ has two distinctive
features that interact with the strict-$h=0$ degeneracy:

\begin{itemize}
\item The price function $P(u)$ maps to $(0, 1)$, the same range as the
state probability.  This makes $\PFR(u) = \sigma(\tau \sum_k u_k)$ a
particularly clean closed form.
\item The Bayes posterior $\mu_k$ is a single number (probability of $v
= 1$), not a distribution.  This makes the slice-integral conditioning
algebraically tractable.
\end{itemize}

Together these features make the Bayes calculation under $P = \PFR$ a
"sufficient-statistic" computation that closes in closed form, giving
$\mu_k = p$ exactly for all $k$.

Under a continuous payoff, neither of these features holds: $P$ would
map to an interval different from the state space, and Bayes would
produce a distribution rather than a number.  The closed-form
calculation $\mu_k = p$ would fail, and $\PFR$ would (generically)
NOT be a fixed point of the strict-$h=0$ operator.

This suggests the GS degeneracy is a knife-edge phenomenon of the
binary payoff case.  The exact extent of this phenomenon under other
payoff distributions is an open research question.

\section{Discussion: implications for empirical work}

The numerical certification has direct implications for empirical
calibration of partial-revelation REE:

\begin{itemize}
\item The asymptotic $\deficit \propto 1/\gamma$ at fixed $\tau$
provides a moment condition relating risk aversion to observed
price-information content.
\item The K-invariance of the deficit means that empirical studies do
not need to know exactly how many traders are active — the deficit
depends primarily on $(\tau, \gamma)$.
\item Trading volume $\propto K$ provides an independent moment for
calibrating the number of traders.
\end{itemize}

\section{Discussion: implications for theory}

The existence of a noise-free partial-revelation mechanism has
implications for the Grossman-Stiglitz literature:

\begin{itemize}
\item The standard noise-trader assumption is not necessary for
partial revelation: risk aversion under binary payoff suffices.
\item The exact role of payoff structure deserves further investigation:
the binary case admits a closed-form revealing equilibrium that
generates the GS degeneracy, but continuous payoff cases may not.
\item The kernel bandwidth $h$ is mathematically equivalent to a
small "noise" parameter (via the convolution structure of the
regularization).  This suggests a deep connection between the
$h$-regularization selection mechanism and the noise-trader selection
mechanism.
\end{itemize}

\section{Conclusion (extended)}

This study has established a partial-revelation equilibrium in a
small-but-canonical setting through extensive numerical and analytic
work.  The 255 certified equilibria at $10^{-15}$ extended-precision
residual constitute the first systematic catalog of CRRA-binary REE
solutions.  The K-invariance finding extends naturally to multiple
crowds.  The Grossman-Stiglitz degeneracy at strict $h = 0$ is the
methodological surprise of the project and reframes the kernel
bandwidth as an equilibrium selection mechanism rather than a
numerical convenience.

Open directions:
\begin{itemize}
\item Continuous payoff variant: does the GS degeneracy persist?
\item Computer-assisted proof: convert Theorem 1's $\chi < 1$ condition
to global-in-$h$ via interval arithmetic at a certified anchor.
\item Multi-asset extensions: the discretization scales as $G^{KJ}$ for
$J$ assets, requiring symmetric framework extensions.
\item Strategic information acquisition: at what level of information
cost does the partial-revelation equilibrium become incentive-compatible
for traders to acquire information?
\end{itemize}





\section{Numerical recipes: code-level walkthroughs}
\label{app:recipes}

For reviewers and reproducers, this appendix walks through the production
solver pipeline at the code level.

\subsection{Recipe 1: solving one cell with the kernel operator}

Input: $(\tau, \gamma)$, target grid $G$, optional warm-start P.

Steps:
\begin{enumerate}
\item Initialize $P_0$ to the no-learning prediction.
\item For $G' \in \{9, 13, 17, G\}$ in increasing order:
\begin{itemize}
\item Build $G'$-grid, halo cells.
\item Initial guess: upsample $P_{G' - 4}$ to $G'$ resolution.
\item scipy.optimize.newton\_krylov with $f_{\rm tol} = 10^{-12}$,
$maxiter = 30$.
\end{itemize}
\item Return $P_G$.
\end{enumerate}

Typical wall: 5--20 seconds per cell at $G=21$.

\subsection{Recipe 2: longdouble polish}

Input: certified P at $G=21$ from Recipe 1.

Steps:
\begin{enumerate}
\item Convert $P$ to numpy.longdouble.
\item Newton-GMRES(40) iteration:
\begin{itemize}
\item Outer Newton iter: build matvec via finite-difference with $\epsilon = 3 \times 10^{-10}$.
\item Inner GMRES with relative tolerance $10^{-3}$, max 40 Krylov vectors.
\item Accept Newton step.
\end{itemize}
\item Stop when $\|F\|_\infty < 10^{-15}$ OR max 6 Newton iterations.
\item Return polished P.
\end{enumerate}

Typical wall: 30--60 seconds per cell at $G=21$.

\subsection{Recipe 3: $\gamma$-continuation rescue}

Input: cell $(\tau, \gamma)$ that stalled in Recipe 1.

Steps:
\begin{enumerate}
\item Find nearest neighbor $\gamma'$ at the same $\tau$ that was accepted.
\item Use that converged $P$ as warm-start for $\gamma$.
\item Run Recipe 1 with the warm-start.
\end{enumerate}

This rescues most $\gamma$-direction stalls.  For $\tau$-direction stalls,
use Recipe 4 below.

\subsection{Recipe 4: $\tau$-continuation rescue}

Input: cell $(\tau, \gamma)$ where Recipe 3 also failed.

Steps:
\begin{enumerate}
\item Find nearest neighbor $\tau' < \tau$ that has an accepted solution
at this $\gamma$.
\item March $\tau \to \tau$ in increments of $0.1$, using the previous
$\tau'$ converged $P$ as warm-start for $\tau' + 0.1$.
\item Final step at the target $\tau$.
\end{enumerate}

This recovers most cells in the high-$\gamma\tau$ region.

\subsection{Recipe 5: deep G-ladder for continuum extrapolation}

Input: certified cell at $G=21$.

Steps:
\begin{enumerate}
\item For $G' \in \{25, 29, 33, 37\}$:
\begin{itemize}
\item Index-upsample $P_{G'-4}$ to $G'$.
\item Newton-Krylov polish at $G'$.
\item If $F > 10^{-8}$: mark as stalled, do not use in extrapolation.
\end{itemize}
\item Continuum fit: q-free, q=1, q=2 fits over nailed rungs.
\item Continuum estimate = q=1 value; error bar = max spread.
\end{enumerate}

\subsection{Recipe 6: certified deficit and welfare metrics}

For each certified cell:
\begin{enumerate}
\item Compute the inner cube $G=21$ array $P$.
\item Compute $\deficit = 1 - R^2$ of linear regression of
$\mathrm{logit}\,P$ on $\sum_k u_k$ over the inner cube.
\item Compute trading volume $\mathrm{TV} = \E\,|x^*(\mu, p)|$
under the true joint measure.
\item Compute the value-of-information decomposition by binning the
joint distribution and computing $\mathrm{CE}_{\mu_{\rm FR}}$,
$\mathrm{CE}_{\rm informed}$, $\mathrm{CE}_{\rm priceonly}$,
$\mathrm{CE}_{\rm prior}$ under the true measure.
\end{enumerate}

\section{Glossary of notation and acronyms}

\begin{itemize}
\item $\tau$: signal precision (inverse signal variance)
\item $\gamma$: CRRA relative risk aversion
\item $K$: number of traders ($K \in \{3, 4, 5\}$ in this study)
\item $G$: discretization resolution (cells per axis)
\item $u_k$: trader $k$'s private signal, centered ($u_k = s_k - 1/2$)
\item $v$: binary asset payoff ($v \in \{0, 1\}$)
\item $P(u_1, \ldots, u_K)$: equilibrium price function
\item $\PFR = \sigma(\tau \sum_k u_k)$: fully revealing equilibrium price
\item $\mu_k$: trader $k$'s posterior of $v=1$ given own signal and price
\item $A_v(p, u_k)$: evidence integral (Equation~\ref{eq:evidence})
\item $\Phi$: the equilibrium operator (Bayes-then-clear)
\item $\Phi_h, \Phi_0$: kernel-smoothed (h>0) and strict (h=0) versions
\item $\deficit = 1 - R^2$: revelation deficit
\item $\Vipub$: value of observing the price vs.\ prior baseline
\item $\Vipriv$: value of own signal above the price
\item $\ViFR$: welfare gap to full revelation
\item $\mathrm{TV}$: trading volume (per agent)
\item REE: rational-expectations equilibrium
\item CRRA: constant relative risk aversion
\item CMM: co-area moving mesh (the strict-h=0 program)
\item GS: Grossman-Stiglitz
\item PR: partially revealing
\item FR: fully revealing
\item LM: Levenberg-Marquardt
\end{itemize}

\section{Final remarks}

This summary documents an extensive computational project on K=3 CRRA REE
with binary payoff and no noise traders.  The economic finding (partial
revelation from risk aversion alone), the numerical certification
(255 cells at $10^{-15}$), the K-invariance result (deficits coincide
within $\pm 25\%$ across $K \in \{3,4,5\}$), and the Grossman-Stiglitz
degeneracy at strict $h=0$ are the main contributions.  The project's
codebase, certified solutions, and exhaustive intermediate logs are all
in the repository for reproducibility and future extensions.





\section{Extended proof of Theorem 1 (full detail)}
\label{app:fullproof}

We give an extended proof of Theorem~\ref{thm:existence}.

\subsection{Setup}

Fix $(\tau, \gamma, U) \in (0, \infty)^3$ and regularization parameters
$(h, \delta) \in (0, \infty)^2$ with $\chi = 12 \tau U / h^2 < 1$.

Define the function class $\mathcal{K}(B, L)$ as in Section 8.  It is
compact convex in sup-norm by the Arzela-Ascoli theorem (continuous
functions on a compact domain, uniformly bounded, uniformly Lipschitz
form a compact set).

Define $\Phi_{h, \delta}: \mathcal{K}(B, L) \to C(Q)$ as the
regularized equilibrium operator.

\subsection{Lemma 1: bounded evidence ratio}

\textbf{Statement.} For $\Pi \in \mathcal{K}(B, L)$, the
regularized log-evidence-ratio satisfies $|b(u, \Pi)| \le 2\tau U$ at
every $u \in Q$.

\textbf{Proof.} The integral defining $b$ involves Gaussian densities
$f_0, f_1$ over the support $|u_j - 1/2| \le U + 1/2$ and $|u_j + 1/2|
\le U + 1/2$.  The log-likelihood ratio $\log(f_1(u_j)/f_0(u_j)) = \tau u_j$
has range $[-\tau U, \tau U]$ over the support.  Summing over $K-1 = 2$
other agents: $|b| \le 2 \tau U$.  $\square$

\subsection{Lemma 2: derivative bounds}

\textbf{Statement.} For $\Pi \in \mathcal{K}(B, L)$,

\begin{align}
\biggl|\frac{\partial b}{\partial \hat\pi}\biggr| &\le \chi = \frac{12 \tau U}{h^2} \\
\biggl|\frac{\partial b}{\partial u_j}\biggr| &\le \beta = \frac{4U}{\delta^2}, \quad j \ne k
\end{align}

\textbf{Proof sketch.} The derivative with respect to the observed price
$\hat\pi$ acts through the kernel $K_h$.  By Gaussian moment estimates,
$|\partial K_h(x)/\partial x| \le K_h(x) \cdot |x|/h^2 \le 1/(h^2 \sqrt{2 \pi e})$.
Multiplying by the range of $|b|$ (Lemma 1) and the integration domain
size $(2U+1)$, we get the $12 \tau U / h^2$ bound after algebraic
simplification.

The derivative with respect to $u_j$ acts through the own-signal
mollifier $\rho_\delta$ with bandwidth $\delta$.  Similar Gaussian
moment estimate gives the $4 U / \delta^2$ bound.  $\square$

\subsection{Lemma 3: CRRA non-expansive aggregation}

\textbf{Statement.} The CRRA-clearing operator on log-odds
$\pi^* = \mathrm{logit}\, p^{\rm clear}(\mu_1, \ldots, \mu_K; \gamma)$
satisfies
\begin{equation}
\frac{\partial \pi^*}{\partial m_k} = w_k \ge 0, \quad \sum_k w_k \le 1.
\end{equation}
The strict inequality $\sum_k w_k < 1$ holds whenever any $\mu_k \ne p$
(i.e.\ whenever some trader disagrees with the price).

\textbf{Proof sketch.} Differentiating the clearing condition $\sum_k
x_k^* = 0$ implicitly in $m_k = \mathrm{logit}\,\mu_k$ gives
$\partial \pi^*/\partial m_k = (\partial x_k^*/\partial m_k) / \sum_j (\partial x_j^*/\partial \pi^*)$.
Both numerator and denominator can be computed in closed form:
\begin{align}
\frac{\partial x_k^*}{\partial m_k} &= \frac{R_k p (1-p) / \gamma}{[(1-p) + R_k p]^2} \\
\sum_j \frac{\partial x_j^*}{\partial \pi^*} &= -\frac{1}{p(1-p)} - \frac{1}{\gamma} \sum_j \frac{R_j p(1-p)}{[(1-p) + R_j p]^2}
\end{align}
Their ratio gives $w_k$.  The sum $\sum_k w_k$ has a clean closed
form involving the "Jensen slack" terms $J_k = (R_k - 1)^2 p (1-p) / D^2$
with $D = (1-p) + R p$:
\begin{equation}
\sum_k w_k = 1 - \frac{(\gamma + p(1-p)) \sum_k J_k}{\sum_k [(R_k - 1) D + R_k p(1-p)/\gamma]}
\end{equation}
The denominator is positive; the numerator is non-negative
($J_k \ge 0$); equality holds iff all $R_k = 1$ (all $\mu_k = p$). $\square$

\subsection{Lemma 4: invariance of $\mathcal{K}(B, L_\star)$}

\textbf{Statement.} Under $\chi < 1$, $\Phi_{h, \delta}(\mathcal{K}(B, L_\star)) \subset \mathcal{K}(B, L_\star)$.

\textbf{Proof.}
For $\Pi \in \mathcal{K}(B, L_\star)$, we need to show that
$\Phi_{h, \delta}\Pi$ also satisfies:
\begin{itemize}
\item permutation-symmetric: trivial (the operator preserves
permutation symmetry of the input).
\item odd: trivial under sign-flip symmetry of the model.
\item bounded $|\Phi \Pi| \le B$: by composition, the output's logit
is the CRRA-cleared log-odds of the trader posteriors, each of which is
$\tau u_k + b$ with $|b| \le 2 \tau U$ (Lemma 1).  The clearing is a
non-expansive (Lemma 3) average of these, so the cleared logit is bounded
by $\tau U + 2 \tau U = 3 \tau U = B$.
\item Lipschitz $L_\star$: by the chain rule, the Lipschitz constant of
$\Phi \Pi$ is the product of:
(i) the Lipschitz constant of the CRRA-clearing operator on log-odds,
$\le 1$ by Lemma 3;
(ii) the Lipschitz constant of the posterior map $\mathrm{logit}\,\mu = \tau u + b$,
which is $\tau + \chi L + \beta$ by Lemmas 1, 2.  Solving the fixed-point
inequality $L_\star = \tau + \chi L_\star + \beta$ gives $L_\star = (\tau + \beta)/(1 - \chi)$. $\square$
\end{itemize}

\subsection{Schauder application}

$\mathcal{K}(B, L_\star)$ is compact convex in sup-norm.  $\Phi_{h, \delta}$
is continuous (in fact smooth) and maps $\mathcal{K}$ to itself (Lemma 4).
By Schauder's fixed-point theorem, $\Phi_{h, \delta}$ has a fixed point
in $\mathcal{K}(B, L_\star)$.

\subsection{Proof of $\Pi^\star \ne \mathrm{logit}\,\PFR$}

If $\chi \le 1/3$ and $\beta < \tau(2 - 3\chi)$, then
$L_\star = (\tau + \beta)/(1 - \chi) < 3 \tau$.  The fully-revealing $\PFR$
has $\mathrm{logit}\,\PFR(u) = \tau \sum_k u_k$, with Lipschitz constant
$3\tau$ in sup-norm.  Since $\Pi^\star \in \mathcal{K}(B, L_\star)$ has
Lipschitz constant $\le L_\star < 3\tau$, $\Pi^\star \ne \mathrm{logit}\,\PFR$. $\square$

\section{Why this matters for empirical asset pricing}

The findings here have potential implications for empirical asset
pricing models that calibrate to observed price-revealing power:

\begin{itemize}
\item The deficit's $\propto 1/\gamma$ asymptote provides a quantitative
mapping from observable price-information content to underlying risk
aversion, independent of any noise-trader assumption.
\item The K-invariance result simplifies calibration: the deficit
depends primarily on $(\tau, \gamma)$, not on the (often unknown)
crowd size.
\item Trading-volume proportionality $\propto K$ provides an
independent moment for crowd size, decoupling the deficit's identification
from the volume's identification.
\end{itemize}

\section{Final summary table: the 7 main findings}

\begin{table}[H]\centering\small
\begin{tabular}{lp{8cm}}
\toprule
Finding & Statement\\
\midrule
1. Existence (regularized) & Theorem 1: regularized PR-REE exists for any $(\tau, \gamma)$
when $\chi < 1$ and the fixed point is strictly disjoint from $\PFR$.\\
\midrule
2. Numerical certification & 255 individual equilibria at extended-precision
residual $\le 10^{-15}$ across 12 $\tau$ values and 23 $\gamma$ values.\\
\midrule
3. Continuum limits & Immortal anchor: $\deficit_\infty = 0.268 \pm 0.010$.
Max-deficit corner: $\deficit_\infty = 0.278 \pm 0.008$.  G=21 values within 3\% of continuum.\\
\midrule
4. Welfare decomposition & $\Vipub \gg \Vipriv > 0$ at every parameterization;
$\ViFR$ vanishes proportionally to $\deficit$.\\
\midrule
5. K-invariance & Deficits coincide within $\pm 25\%$ per cell across $K \in \{3,4,5\}$;
slope of $\mathrm{logit}\,P$ on $\sum_k u_k$ equals $\tau/K$ to three sig figs everywhere.\\
\midrule
6. GS degeneracy & $\PFR$ is the unique strict-$h{=}0$ fixed point under binary payoff,
certified to $10^{-15}$ in 9 test cells.\\
\midrule
7. Implication & Kernel bandwidth $h > 0$ is the equilibrium selection mechanism for the
nontrivial PR branch; no noise traders required.\\
\bottomrule
\end{tabular}
\caption{Final summary of the seven main findings.}
\end{table}



\input{cells_table}\input{k_table}\input{vi_table}

\end{document}
